% =====================================================================
%  P4.  Cell-mean fluctuations: exact Rademacher-null variances for cell
%       means of arithmetic fields, with an application to
%       sum_{n<N} Lambda(n) mu(N-n).
%
%  Compiles standalone with pdflatex, run twice.
% =====================================================================
\documentclass[11pt,a4paper]{article}

\usepackage{amsmath}
\usepackage{amssymb}
\usepackage{amsthm}
\usepackage[margin=1in]{geometry}
\usepackage[colorlinks=true,linkcolor=blue,citecolor=blue]{hyperref}

\theoremstyle{plain}
\newtheorem{theorem}{Theorem}
\newtheorem{corollary}[theorem]{Corollary}
\newtheorem{proposition}[theorem]{Proposition}
\newtheorem{lemma}[theorem]{Lemma}
\theoremstyle{definition}
\newtheorem{observation}[theorem]{Observation}
\newtheorem{measurement}[theorem]{Measurement}
\theoremstyle{remark}
\newtheorem{note}[theorem]{Note}

\renewcommand{\SS}{\mathfrak{S}}
\newcommand{\AAA}{\mathfrak{A}}

% typographic slack: let TeX stretch a line rather than run into the margin
\emergencystretch=3em

\title{Cell-mean fluctuations in an arithmetic convolution field:\\
exact Rademacher-null variances and a permutation control}
\author{Youngrok Lee\\
\small Independent researcher, Republic of Korea}
\date{}

\begin{document}
\maketitle

\begin{abstract}
Let $Z(N)=C(N)/\sqrt{V(N)}$ with
$C(N)=\sum_{n<N}\Lambda(n)\mu(N-n)$, and partition a band of even $N$
into cells indexed by how many of $3,5,7,11,13$ divide $N$. Whether a
cell mean departs from the band mean needs a reference scale, and the
count-based one misses the null scale here by a factor that grows with
the band.

Under the null of independent Rademacher signs for the $\mu(v)$ on
their support, the variance of any cell contrast is a
doubly centred quadratic form in the covariance kernel and reduces to
three convolutions: the exact Rademacher-null scale is written down,
not estimated. A uniformly random cell of the same size recovers the
finite-population formula exactly; the divisibility cells do not.
Measured over eight octaves to $1.6\cdot10^{7}$, the exact null
standard deviation shrinks by $1.21$ where a count predicts $11.3$,
and a label permutation preserving every cell size collapses it by
factors $3.8$ to $105$.

For its size we offer a descriptor: the excess $D_c$ of the shift's
singular series over same-cell pairs, whose main term is rational with
no fitted parameter. That the residue structure determines $D_c$ was
committed for a second, differently shaped partition before it was
measured; the exact measurement falls inside the registered interval
at every cell, and that transfer is the evidence offered. That $D_c$
in turn determines the null floor rests on a correlation of $0.9805$
across six cells, computed after the fact and not transferred. Every
registered rule is reported with its verdict, refuted ones included.
No distributional law for $Z$ and no Goldbach consequence is asserted,
and the decay of the effect with $N$ is not determined here.
\end{abstract}

% =====================================================================
\section{Introduction}

\subsection{The setting}

Let $N$ range over the even integers of a band, let
\begin{equation}\label{eq:field}
  C(N)=\sum_{n<N}\Lambda(n)\mu(N-n),
  \qquad
  V(N)=\sum_{v<N}\mu^2(v)\Lambda(N-v)^2,
  \qquad
  Z(N)=\frac{C(N)}{\sqrt{V(N)}},
\end{equation}
so that $Z$ has second moment $1$ under the null in which the $\mu(v)$
on their own support are replaced by independent signs; the exact
asymptotic $V(N)\sim\prod_{q\nmid N}(1-1/(q(q-1)))\,N\log N$ is
Proposition~\ref{prop:V} in Appendix~\ref{sec:scale} and is quoted for
scale only: no computation below uses it, every one dividing by $V(N)$
as computed. The heuristic of Observation~\ref{obs:coh} does read the
scale $N\log N$ off it, and nothing else here does. The field \eqref{eq:field} is the scalar to
which the Huang--Li reduction of binary Goldbach reduces: their
Theorem~1, with Corollary~1 drawing the Goldbach conclusion from it,
turns on the equidistribution of $\Lambda(n)\mu(N-n)$ in arithmetic
progressions \cite[Thm.~1, Cor.~1]{HL}. That is imported, not derived
here, and used for nothing below, since
\S\ref{sec:notclaimed} refuses the Goldbach consequence outright.
Partition the
band into \emph{cells} indexed by \emph{depth} $d$, the number of
$3,5,7,11,13$ dividing $N$. The question is whether
$m_c-\overline m$, the difference between the cell mean of $Z$ and the
band mean, is larger than chance.

Answering it requires a reference scale for $m_c-\overline m$. The
reflex is $\mathrm{sd}(Z)/\sqrt{n_c}$ with $n_c$ the cell count. That
the reflex fails for correlated observations is standard --- the
variance of a weighted mean is $w^{\top}\Sigma w$, and the ratio to
the independent formula is the design effect \cite{Kish}. What is not
standard, and is what this paper supplies, is that for the field
\eqref{eq:field} the covariance $\Sigma$ is \emph{exactly computable}
under a sign null; that cells cut by divisibility align with it; that
a uniformly random cell of the same size recovers the
finite-population formula exactly; and that the size of the effect is
predicted by an excess of the shift's singular series over same-cell
pairs.

\subsection{Results}

\begin{itemize}
\item Lemma~\ref{lem:cellmom} gives $\mathrm{Var}(m_c-\overline m)$
exactly under the sign null, as three convolution terms.
\item Observation~\ref{obs:coh} takes up its decay: over the eight
measured octaves the shallow cells are consistent with
$(\log N)^{-1/2}$ and not with $n_c^{-1/2}$. The first is a heuristic
read against a finite range, not a proved rate; what is established is
that the count-based rate is wrong.
\item Lemma~\ref{lem:placebo} is the control that separates a
property of the cell \emph{sizes} from a property of the
correspondence between cells and arithmetic.
\item Measurement~\ref{meas:placebo} runs it, and finds that the
effect survives while the floor itself does not --- so the floor is a
second measurement of the same correspondence.
\item Proposition~\ref{prop:scaleinv} shows the quantity predicting
the floor's size is asymptotically scale-free: as the band grows it
converges to a constant fixed by the cell's residue classes, so its
one-step exponent tends to zero wherever that constant is nonzero.
\item Measurement~\ref{meas:maincoef} computes that quantity's main
term in exact rational arithmetic and finds it reproduces the measured
profile at all six depths with no fitted parameter --- which both
supplies the hypothesis the proposition needs (that the main term does
not vanish) and turns the proposition from a statement that the size
does not drift into one that names it.
\item Measurement~\ref{meas:secondcell} carries the same computation to
a \emph{second} cell shape, predicting its profile before measuring it
and measuring it exactly. The prediction holds where the rules could
score it, and the small
correction it carries decomposes with no free parameter into a divisor
deficit and a covariance --- the second of which is exactly zero at the
only cell where the first was ever calibrated.
\end{itemize}

\subsection{What is not claimed}\label{sec:notclaimed}

\begin{itemize}
\item \textbf{No decay exponent for the effect is quoted, and the
shallow depths are not uniformly quiet.} That the amplitude at the
three shallowest depths clears nothing at any measured scale cannot be
read off two octaves, and the two whose per-depth amplitudes were
printed are not the quietest two: one is the smallest of the eight and
the other the third smallest.
The run behind Measurement~\ref{meas:notzero} prints all eight, and at
depth~$1$ the amplitude runs
\[
  +2.71,\ +2.18,\ +1.96,\ +1.72,\ +1.34,\ +1.11,\ +0.87,\ +0.69
\]
from $(6.25\cdot10^4,1.25\cdot10^5]$ up to
$(8\cdot10^6,1.6\cdot10^7]$ --- the value at the smallest octave
exceeding the $-1.55$ that depth~$3$ carries at the largest, where the
paper calls the deep cells the signal. So no threshold separating
``clears'' from ``does not'' is asserted at the shallow depths. What
is asserted is the shape: at each of depths $0$, $1$ and $2$ the
amplitude falls monotonically across all eight octaves, and no
exponent is fitted to that fall. Whether the exponents differ by depth
is not decided by the accessible range.
\item \textbf{No distributional law for $Z$ is asserted.}
\item \textbf{No consequence for the Goldbach problem.} Nothing
measured here bounds $C(N)$ in either direction. No sweep over
parameterisations is carried out, and none is claimed.
\item \textbf{That $D_c$ predicts the floor's size is measured, not
derived.} The variance and $D_c$ are the same doubly centred mean of
two kernels, and a constant ratio between them would follow if those
two kernels differed across the band by a factor the centring leaves
alone. No such statement is proved here. What is measured is the
ratio: its full width is $25.4\%$ of its mean, or $19.8\%$ once
$\mathrm{se}_c^2$ is matched to $D_c$'s convention term by term ---
matched that way the row is still monotone from depth~$1$ on, with a
last step nine times smaller than the printed one --- while
$D_c$ moves by a factor $68$, and the correlation $0.9805$ has an exact enumeration null over
all $720$ relabellings. That is one octave read across depths.
Read across scale instead, the same ratio drifts by about half as much
again (Measurement~\ref{meas:Dc}), so what is measured is a depth profile at
a fixed scale and $D_c$ is not all of what fixes the size.
\eqref{eq:JKJ} is a difference of three means and not of two, and
Measurement~\ref{meas:Dc} records what dropping the third would cost.
\item \textbf{Observation~\ref{obs:coh} is offered for its form, not
its constant.} Its heuristic is quantitatively loose by factors
running from $2.49$ at depth $0$ to $6.20$ at depth $5$ --- the $2.5$
is the shallowest cell, not the range, whose lower end is the $2.45$
at depth~$1$; the exponent is confirmed
directly instead.

\item \textbf{The statements here are not equally checked.} The
derivation of Observation~\ref{obs:coh} and the statement of
Lemma~\ref{lem:placebo} have had less scrutiny than the rest, and
Proposition~\ref{prop:scaleinv} has had no reading by a subject
expert. What was checked by what is set out in Supplement~S1
(Appendix~\ref{sec:repro}).
\end{itemize}

% =====================================================================
\section{The floor, in closed form}\label{sec:floor}

The identity this section rests on is not arithmetic. It is the
statement that a cell contrast is a linear functional of the observation
vector, so its variance is the quadratic form $w^{\!\top}\Sigma w$ in
whatever covariance $\Sigma$ the observations carry. We state that first,
in the generality in which it is true, and specialise afterwards.

\begin{lemma}[contrast variance under an arbitrary kernel]\label{lem:contrast}
Let $X=(X_N)_{N\in a}$ be any square-integrable random vector indexed by
a finite set $a$ of size $n$, with covariance
$\Sigma_{NN'}=\mathrm{Cov}(X_N,X_{N'})$. For $c\subseteq a$ of size
$n_c\ge1$ let $m_c=n_c^{-1}\sum_{N\in c}X_N$ and
$\overline m=m_a$. Writing
$Q_{cd}=\sum_{N\in c}\sum_{N'\in d}\Sigma_{NN'}$,
\begin{equation}\label{eq:contrast}
  \mathrm{Var}\bigl(m_c-\overline m\bigr)
  \;=\;\frac{Q_{cc}}{n_c^{2}}
  \;-\;\frac{2\,Q_{ca}}{n_c\,n}
  \;+\;\frac{Q_{aa}}{n^{2}}.
\end{equation}
\end{lemma}

\begin{proof}
$m_c-\overline m=w^{\!\top}X$ with $w_N=n_c^{-1}\mathbf 1_{N\in c}-n^{-1}$,
so the variance is $w^{\!\top}\Sigma w$. Expanding gives the four
products of the two terms of $w$ against themselves; they collapse to
three sums because the two cross terms are equal, $Q_{ac}=Q_{ca}$ by
symmetry of $\Sigma$.
\end{proof}

Nothing above uses independence, stationarity, or a distribution. What
the null supplies is not the identity but the kernel: it makes $\Sigma$
diagonal in a basis where it can be summed, and hence makes $Q_{cd}$
computable rather than merely defined. That is the content of the next
lemma, and it is where the arithmetic enters.

\begin{lemma}[exact cell moments]\label{lem:cellmom}
Let $c$ be any set of $n_c$ of the even $N$ of a band of size $n$ ---
the proof uses nothing beyond $c\subseteq a$, and in this paper the
cells do partition the band --- let $m_c$ and $\overline m$ be
the means of $Z$ over $c$ and over the band, and set
\[
  u_c(v)\;=\;\sum_{N\in c}\frac{\Lambda(N-v)}{\sqrt{V(N)}},
  \qquad
  Q_{cd}\;=\;\sum_v\mu^2(v)\,u_c(v)\,u_d(v).
\]
Then, when the $\mu(v)$ on the surviving support are replaced by
independent signs,
\begin{equation}\label{eq:cellmom}
  \mathrm{Var}\bigl(m_c-\overline m\bigr)
  \;=\;\frac{Q_{cc}}{n_c^{2}}
  \;-\;\frac{2\,Q_{ca}}{n_c\,n}
  \;+\;\frac{Q_{aa}}{n^{2}},
\end{equation}
where $a$ denotes the whole band. Every term is one convolution of two
explicit finite sequences: the identity is exact and no simulation is
involved. Its numerical evaluation is a separate matter --- the
convolutions below are carried out in double precision
(Appendix~\ref{sec:repro}), so each printed floor is this exact formula
evaluated to floating-point accuracy and not an exact rational.
\end{lemma}

\begin{proof}
Under independent signs the covariance carries $\mu^4$, and $\mu^4=\mu^2$
because $\mu$ takes only $0$ and $\pm1$; so
$E[\mu^2(v)\varepsilon(v)\cdot\mu^2(v')\varepsilon(v')]
 =\mu^2(v)\delta_{vv'}$ and
$\mathrm{Cov}(n_cm_c,n_dm_d)=Q_{cd}$ exactly, which is the $Q_{cd}$ of
Lemma~\ref{lem:contrast} for this $\Sigma$. That lemma then gives the
three terms.
\end{proof}

Throughout, the \emph{null standard deviation} of a cell contrast means
the square root of the quantity this lemma computes: the exact standard
deviation of $m_c-\overline m$ under the Rademacher null, at the
observed cell sizes and the observed support of $\mu^2$. It is a scale,
not a tail bound. The lemma fixes it exactly, but no distributional law
is claimed for $m_c-\overline m$, so no probability is attached to any
multiple of it and none is quoted below.

\begin{note}[all three terms are needed]\label{note:threeterms}
The substitution $u_c(v)\approx n_c/\sqrt V$, which gives the right
scale for $Q_{cc}/n_c^2$, is not specific to $c$; applied to all three
terms it returns the same value $S$ for each and hence $S-2S+S=0$. So
it cannot be used to estimate the variance of the difference. Exactly,
the variance is smaller than $Q_{cc}/n_c^2$ by a factor that is
constant \emph{in $N$} rather than noise --- constant for a fixed
cell, not across cells: at one $N$ it reads $0.113$, $0.016$, $0.137$,
$0.289$, $0.427$, $0.552$ from depth~$0$ to depth~$5$. That row is not
monotone either --- its minimum $0.016$ sits at depth~$1$, a seventh of
the value at depth~$0$ --- so the spread across cells is a factor $34$
and not the factor $5$ that reading the two ends alone would give.
Across the eight octaves it holds to four
digits at depths $0$ and $1$ and to three at depths $2$ and $3$,
drifts by $2.2\%$ at depth~$4$ and by
$30\%$ at depth~$5$, whose cell grows from $2$ members at the lowest
octave to $266$ at the highest: the
constancy is a statement about the shallow cells, where there are
enough $N$ for it to be one.
\end{note}

\begin{measurement}[the closed form, checked against simulation]%
\label{meas:mc}
The lemma does not need simulation, but simulation is the only thing
that would catch an error in the formula, and every normalised cell
mean below is divided by it. Run in the band $(10^5,2\cdot10^5]$ with
$2000$ draws:
\[
\begin{array}{r|cccccc}
 \text{depth} & 0 & 1 & 2 & 3 & 4 & 5\\
 n_c & 19181 & 21097 & 8183 & 1421 & 115 & 3\\\hline
 \text{closed form} & 0.140080 & 0.052724 & 0.167579 & 0.273630
   & 0.378073 & 0.639330\\
 \text{MC}/\text{closed} & 0.9920 & 0.9993
   & 0.9918 & 1.0044 & 1.0157 & 1.0320
\end{array}
\]
--- worst deviation $0.0320$, at the depth whose cell holds three
elements, and the six ratios straddle $1$ three and three.

The six ratios are estimated from the \emph{same} draws, so their
deviations are strongly correlated and their common sign would be one
observation, not six. At $2000$ draws the Monte-Carlo precision is
$1/\sqrt{2\cdot1999}=0.015815$, and the six deviations stand to it
in $0.507,\,0.041,\,0.519,\,0.275,\,0.992,\,2.025$. Against that bar
only the cell of three elements clears, and clearing it is not
agreement.

\textbf{That bar is a convention and not the precision.}
$1/\sqrt{2(n-1)}$ is the Gaussian value of a delta-method
approximation to the relative standard error of a sample standard
deviation; it is not exact for any distribution, and the sampling
variance of a sample standard deviation has no closed form here. What
is exact is one level down: for the sample \emph{variance} of $n$
independent draws the sampling variance is
$n^{-1}(\mu_4-\tfrac{n-3}{n-1}\sigma^4)$, and $\mu_4$ is available,
since each contrast is a weighted Rademacher sum
$\sum_v a_v\varepsilon_v$ with
$\mathbb E S^4=3(\sum a^2)^2-2\sum a^4$.

Two things follow and only the second is used. The kurtosis is
$3-2\sum a^4/(\sum a^2)^2$, below the Gaussian $3$ at every cell.
So the approximation the bar comes from is, to leading order, an
\emph{over}estimate, and it is one that varies from cell to cell since
the weights do. The direction is what matters for the sentence above:
a cell-specific bar is tighter, so ``only the cell of three clears''
is a statement about the common bar and not about the cells. What the
common bar supports is the weaker reading, that the deviations are
small on the scale the draw count alone fixes. The cell-wise fourth
moments are not computed here and no exact interval is claimed. That is one deviation
of that size among six correlated estimates and it stays inside the
registered cap of $0.05$, so it refutes nothing; what it shows is that
the resolution and the disagreement sit at the same cell. Of the
remaining five, four sit at $0.52$ of the bar or less, and depth~$4$
sits at $0.992$ of it --- short by $0.8$ per cent of the bar, at the
precision and not below it. At four cells the table resolves nothing,
at one it sits at the bar, and at a number of draws small enough for
the precision to exceed every deviation in the table it would resolve
nothing at all. Two of the four rules
registered in the result file score a $60$-draw quotation this
paper does not make; the ratios of this paragraph are checked in a
block marked post hoc there, and were not registered in advance.

The structural constant of Note~\ref{note:threeterms} is visible
directly: at the top octave of $1.6\cdot10^7$,
$\mathrm{Var}/(Q_{cc}/n_c^2)$ is
$0.113119,\,0.016302,\,0.288567,\,0.551516$ at depths $0,1,3,5$, and
at a quarter of that $N$ the three shallow ratios agree to five
decimals --- $0.113118,\,0.016303,\,0.288571$, the largest gap being
$4\cdot10^{-6}$. Depth $5$ does not: it reads
$0.557654$ there against $0.551516$ at the top, differing in the third
digit, and it is much the smallest cell --- $266$ values of $N$ in the
top octave against $1\,687\,911$ at depth $1$
(\texttt{lab\_cellmom\_montecarlo.py} for the table,
\texttt{lab\_cell\_floor.py} for this paragraph).
\end{measurement}

% =====================================================================
\section{The floor's uncertainty does not fall like a count}%
\label{sec:coh}

\begin{observation}[coherent sums]\label{obs:coh}
The null standard deviation of Lemma~\ref{lem:cellmom} does \emph{not} decay like
$n_c^{-1/2}$. For fixed $v$, about $n_c/\log N$ of the terms of
$u_c(v)$ are nonzero, each of size $\log N/\sqrt V$, so
$u_c(v)\approx n_c/\sqrt V$ and
\[
  \frac{Q_{cc}}{n_c^2}
  \;\approx\;\frac{\sum_v\mu^2(v)}{V}
  \;\asymp\;\frac{1}{\log N},
\]
independently of $n_c$; and by Note~\ref{note:threeterms} the variance
itself is a fixed fraction of this. So the standard error falls
like $(\log N)^{-1/2}$ --- at the shallow cells, which is where
Note~\ref{note:threeterms} says that fraction is constant. It
does not extend to the deepest: the same note records a thirty per
cent drift of the fraction at depth~$5$, and the exponent fitted there
--- $0.107306$, in the ungated rows of \texttt{lab\_cell\_floor.txt}
rather than in Measurement~\ref{meas:exponent}, which prints the three
gated depths --- is the largest of the six. That drift is not the
whole of the break. Separately from it, and not confined to the
smallest cell, $Q_{cc}/n_c^2$ does not follow $1\!/\!\log N$ at any
depth. Over the eight octaves
$Q_{cc}n_c^{-2}\log N$ would be constant if the form held. It falls,
and its value at the lowest octave stands to its value at the highest
in $1.031$ at each of depths~$0$ to~$3$, $1.049$ at depth~$4$ and
$1.621$ at depth~$5$. \textbf{So the form is loose at every depth, and
by about the same amount at the shallow five} --- but not at the
deepest, where the drift is larger by an order of magnitude:
$\log 1.621/\log 1.031$ is $15.8$. The uniformity claimed here is over
depths~$0$ to~$4$ and no further, and depth~$5$ is the single-class
cell that the rest of this section treats separately. The drift is quoted as that ratio and not as a
percentage because a percentage needs a denominator: the two obvious
choices differ by the ratio itself, which is immaterial at $1.03$ and
is not at $1.62$. For the same reason no value is quoted for the
fitted exponent $g$ of that quantity. Its
\emph{sign} is negative at all six depths under
each of the three octave abscissae --- the low endpoint, the midpoint,
the high. Its \emph{value} moves between those three by $0.003721$ at
every one of the six depths, the same number at all of them: the
slope of $\log(Q_{cc}n_c^{-2})$ is untouched, since the abscissae
differ by the constant $\log 2$, so the whole dependence is the shared
$\log N$ factor and an additive term is what a constant swing looks
like. Against $|g|$ that swing is most of the exponent where $|g|$ is
small and four per cent at depth~$5$, which is why no value is
quoted.

That drift is the measured exponent in another coordinate and not a
second reading of it. Where the ratio $\mathrm{Var}/(Q_{cc}n_c^{-2})$
is constant in $N$ --- its spread across the eight octaves is $0$ at
depths~$0$ and~$1$ --- one has $-2b$ for the slope of
$\log(Q_{cc}n_c^{-2})$ in $\log N$, while $g$ is that slope plus
$1/\log N$; so $b=1/(2\log N)-g/2$ identically, and the two sides
cannot disagree.
Nothing in this paragraph should be read as the drift confirming the
exponent.

What the residual of that identity does read is the constancy of the
ratio itself, since that is its only hypothesis, and it is there that
the deepest cells separate. Half the negated slope is $0.039408$
against the fitted $b=0.039388$ at depth~$0$, a difference of
$2\cdot10^{-5}$ where the ratio does not move at all; at depth~$4$ the
difference is $0.0019$ against a ratio spread of $0.0094$, and at
depth~$5$ it is $0.026$ against a spread of $0.165$. So the break at
the deepest cell is the failure of the ratio to hold fixed. The
diagonal $1/n_c$ of $Q_{cc}n_c^{-2}$ rises across the same two
depths --- its share is at most $0.42$ per cent at depths~$0$ to~$3$
and $70.5$ per cent at depth~$5$, at the lowest octave --- and this
measurement does not separate the two, so the diagonal is reported
alongside the break and is not offered as its cause. The exponent is confirmed directly at all six depths and
that confirmation, not the form, is what the paper rests on.
\end{observation}

The heuristic is quantitatively loose and is offered for the
\emph{form} and not the constant: at $N=1.6\cdot10^7$,
\[
\begin{array}{r|cccccc}
 \text{depth} & 0 & 1 & 2 & 3 & 4 & 5\\\hline
 Q_{cc}/n_c^2 & 0.123577 & 0.121208 & 0.146026 & 0.184159 & 0.235171
   & 0.307074
\end{array}
\]
against the heuristic's
$\bigl(\sum_{v<N}\mu^2(v)\bigr)/V(N)=0.049540$
(\texttt{lab\_cell\_floor.py}). The row is not monotone in the depth;
its minimum is $0.121208$ at depth $1$.

\begin{measurement}[the exponent, fitted]\label{meas:exponent}
The form is confirmed directly. Fitting $\mathrm{se}\propto N^{-b}$ to
the exact floor across eight octaves gives
$b=0.039451,\,0.039671,\,0.039388$ at depths $2,1,0$, against the
$0.5$ that a count would give and against $0.036038$, the value of
$1/(2\langle\log N\rangle)$ with $\langle\log N\rangle$ the mean of
$\log N$ over the eight octave midpoints. Nothing turns on the third
decimal: the point is that $0.039$ is near $0.036$ and nowhere near
$0.5$ (\texttt{lab\_cell\_floor.py}).
\end{measurement}

\begin{note}[a count is not the null standard deviation]\label{note:notanerrorbar}
The consequence is not confined to this field. The fit runs over the
eight octaves from $(6.25\cdot10^4,\,1.25\cdot10^5]$ to
$(8\cdot10^6,\,1.6\cdot10^7]$, and the abscissa is the octave midpoint
--- which is what makes $1/(2\langle\log N\rangle)=0.036038$ come out,
with $\langle\log N\rangle=13.8744$. The factor in $N$ is therefore
$1.2\cdot10^7/93750=128$. Over it, $n_c^{-1/2}$ says the null
standard deviation shrinks by $\sqrt{128}=11.3$; it shrinks by $128^{0.0395}=1.21$,
against the $1.19$ that $(\log N)^{-1/2}$ predicts over the same
factor. The result file prints the shrink itself only from depth~$3$
on; at the three gated depths it is $128^b$ from the exponents printed
there, which is the one step taken here. \textbf{An interval built
from a count is therefore, against the exact Rademacher-null scale, a
factor $11.3/1.21=9.3$ too narrow at the top of this range relative to
the bottom}, and in absolute terms the discrepancy is larger still:
at the top octave the exact floor exceeds
$\mathrm{sd}(Z)/\sqrt{n_c}$ by factors of $7.3$ to $158$, with
$\mathrm{sd}(Z)=0.924687$ over that octave --- and $0.929384$ over the
whole fitted field, so the reading of $\mathrm{sd}$ barely moves the
figures. The factor does \emph{not} grow with the cell. What moves with
the cell is the floor itself, which reads $0.1182$, $0.0445$, $0.1414$,
$0.2305$, $0.3169$, $0.4115$ from depth $0$ to depth $5$: it is not
monotone --- depth~$1$ sits below depth~$0$, the same dip that
Measurement~\ref{meas:Dc} reads in $D_c$ and in $\mathrm{se}_c$ --- and
it grows from depth~$1$ onward. That movement is measured and not
predicted by
$Q_{cc}/n_c^2\asymp1/\log N$: that relation is a statement about $N$,
carries no dependence on the cell, and so says nothing about how the
floor moves across depths at one $N$. Section~\ref{sec:notclaimed}
records the same thing from the other side --- the heuristic behind it
is loose by a factor that is $2.49$ at depth~$0$ and rises to $6.20$ at
depth~$5$, with its minimum $2.45$ at depth~$1$ rather than at either
end. The factor
also carries $\sqrt{n_c}$, and $n_c$ runs $1\,534\,466$,
$1\,687\,911$, $654\,281$, $114\,017$, $9\,059$, $266$ across those
depths, which is the larger movement. That row is not monotone either:
its maximum is at depth~$1$, so the largest cell is the one where
Measurement~\ref{meas:Dc} reads the smallest $D_c$ and the smallest
floor. By depth the
factor reads $158$, $62$, $124$, $84$, $33$, $7$
(\texttt{audit\_cellfloor\_sdz.py}). The cause
is that $u_c(v)$ is a coherent sum, not a self-averaging one, and that
the cell is correlated with the arithmetic it sums. Coherence alone
does not give it. A cell drawn at random from the same band is
coherent in the same sense --- every $N$ reads the same $\mu$ and
$\Lambda$ at the same $v$ --- but there $u_c(v)\approx(n_c/n)u_a(v)$
and the three terms of \eqref{eq:cellmom} cancel to leading order, so
Lemma~\ref{lem:permfloor} returns exactly $n_c^{-1/2}$, up to the
finite-population factor, for the average over the permutation. What survives the cancellation is the
departure from that proportionality, and only a cell correlated with
the arithmetic produces one.
\end{note}

% =====================================================================
\section{The control, and what running it shows}\label{sec:placebo}

\begin{lemma}[the placebo key]\label{lem:placebo}
Let the cells be indexed by a labelling $\ell(N)$, and let $\pi$ be a
uniformly random permutation of the even integers in the band. This is
a randomisation \emph{conditional on the observed cell counts}: the
labels are not drawn independently, since the counts are held fixed,
and the field $Z$ is held fixed with them --- the standard
randomisation-inference setting \cite{LR}. Call a
statistic \emph{cell-indexed} when it is a function of the pair
$(\ell,Z)$ --- of which $N$ carry which label, and of the field.
Replacing $\ell$ by
$\ell\circ\pi$ preserves every cell size and leaves the field $Z$
pointwise unchanged, while destroying the correspondence between cells
and arithmetic. Any cell-indexed statistic determined by the multiset
$\{Z(N)\}$ together with the cell sizes --- and so not by that
correspondence --- has its distribution unchanged by the replacement,
and is in fact constant under it. The converse fails, and the lemma
does not assert it: $T(\ell,Z)=Z(N_0)$ at a fixed $N_0$ of the band is
cell-indexed and invariant, since it does not read $\ell$ at all, yet
the multiset and the sizes do not determine it. What the control gives
is the direction stated, used below in contrapositive.
The determining data need not be the sizes alone:
$s^2/n_c$, with $s^2=n^{-1}\sum_N(Z(N)-\overline m)^2$ the variance of
the field over the band, is invariant under the replacement and reads
the field through $s^2$ --- and it is, up to a finite-population
factor, the permutation null \eqref{eq:permnum} that
Lemma~\ref{lem:permfloor} will give. Not the permuted floor
\eqref{eq:permfloor} itself: that is the sign-null expectation of this
one, carrying $E_\varepsilon s^2=1-Q_{aa}/n^2$ where this carries
$s^2$; Note~\ref{note:collapse} reads both off the measured octave. Nor
is the floor at a single $\pi$ determined by the two invariants:
\eqref{eq:JKJ} reads $K$ on the subset actually drawn and $K$ is not
constant across the band. What the invariants fix is the average over
$\pi$, which is what \eqref{eq:permfloor} is --- and the control
reports the ten-draw standard error of that average beside it.

Conversely the replacement supplies a null. The law of a cell-indexed
statistic under $\ell\circ\pi$ is its law under labels drawn
independently of $Z$; and a statistic determined by the two
invariants is constant under the replacement, so its null is a point
mass. Hence a statistic whose true-labelling value that null does not
reproduce is not determined by the two multisets, and what it reads is
the pairing between them --- the correspondence the replacement
destroys. That is the whole of what is claimed: which statistics are
uninformative, not that the permutation law is the right null for any
other question.
\end{lemma}

\begin{proof}
The permutation acts on labels only; the multiset
$\{Z(N):N\ \text{in the band}\}$ and the multiset of cell sizes are
both invariant, and the joint distribution of $(\ell\circ\pi,Z)$ is
that of independent labels. That observation gives both halves: the
first reads it as the law the statistic must have if it sees only
those two invariants, the second as the reference law against which
the true labelling is read. For the second, a statistic determined by the two
invariant multisets takes one and the same value at every $\pi$.
\end{proof}

\begin{measurement}[the cell means are not zero]\label{meas:notzero}
Against the floor of Lemma~\ref{lem:cellmom}, computed from its closed
form, the cell means of the field \eqref{eq:field} are not zero. Over
the eight octaves from $(6.25\cdot10^4,\,1.25\cdot10^5]$ to
$(8\cdot10^6,\,1.6\cdot10^7]$, $\max_c|z_c|$ reads
\[
  10.099,\ 10.954,\ 11.179,\ 12.968,\ 11.885,\ 11.026,\ 10.019,\ 9.064,
\]
No multiple-comparison threshold is asserted for these, and none can
be: the packet carries no family size, no threshold and no null law
for such a statement, and \S\ref{sec:notclaimed} declines the
distributional law any threshold would need. What the numbers are is a maximum over
cells of a standardised deviation against the closed-form null floor,
reported as such. The signal is carried by the deep cells:
at the top octave depths $3,4,5$ sit at $z=-1.55,\,-4.46,\,-9.06$
while depths $0,1,2$ sit at $+0.04,\,+0.69,\,-0.06$
(\texttt{lab\_cell\_floor.py}).
\end{measurement}

\begin{measurement}[the placebo destroys the effect, and the floor
with it]\label{meas:placebo}
Run on the octave $(2\cdot10^6,\,4\cdot10^6]$, with cell sizes
preserved exactly and the floor recomputed from scratch for each
permutation:
\[
\begin{array}{r|cccccc}
 \text{depth} & 0 & 1 & 2 & 3 & 4 & 5\\
 n_c & 383617 & 421978 & 163568 & 28507 & 2263 & 67\\\hline
 z_c\ \text{(true)} & +0.1069 & +1.1133 & -0.2314 & -2.3990
   & -5.9997 & -11.0258
\end{array}
\]
against $\max_c|z_c|$ over ten uniformly random label permutations
(\texttt{default\_rng} seed $20260808$) reading
\[
  0.9359,\ 1.6073,\ 1.9921,\ 1.5687,\ 2.1392,\
  1.1216,\ 1.7178,\ 1.1348,\ 3.2040,\ 1.3192,
\]
mean $1.6741$, never above $3.3$. Ten is a sample and not an
enumeration, and the true value ranking first among the eleven gives a
Monte-Carlo $p$-value of $1/11$, which is the smallest this many
permutations can produce. It is not a bound on the exact permutation
$p$-value, which ten draws cannot bound from above at all; the
enumeration of Measurement~\ref{meas:Dc} below is the
stronger instrument, and this one's limit is the sample's. The
correlation between depth and $z_c$ is $-0.9106$ for the true
labelling and $+0.0042$ on average under permutation. The null does
not reproduce either true value, so by the second half of
Lemma~\ref{lem:placebo} what is detected is the correspondence between
cells and divisibility, not the cell sizes.

The same run corrects a natural reading of Lemma~\ref{lem:cellmom}.
One expects the floor $\mathrm{se}_c$ to be a property of the cell
sizes, hence nearly invariant under permutation. It is not: the floor
\textbf{collapses}, by factors from $3.8$ at depth $5$ to $105$ at
depth $0$ --- $1.2400\cdot10^{-1}$ against $1.1794\cdot10^{-3}$ there
(\texttt{lab\_mask\_placebo.py}).
\end{measurement}

\begin{note}[the floor is itself signal]\label{note:floorsignal}
The mechanism of that collapse is the three-term structure of
\eqref{eq:cellmom}. For a random subset of the band,
$u_c(v)\approx(n_c/n)\,u_a(v)$, and then
$Q_{cc}/n_c^2-2Q_{ca}/(n_cn)+Q_{aa}/n^2$ cancels to leading order,
leaving only the fluctuation around proportionality. An arithmetic
cell breaks that proportionality, and the floor is what is left over.
\textbf{The exact floor is therefore not a noise level but a second
measurement of the same correspondence.} It is not, however, a
conservative version of some truer noise level. Lemma~\ref{lem:cellmom}
is an identity: under the sign null the variance of the arithmetic
cell's numerator \emph{is} the exact floor, so $z_c=-11$ at depth~$5$
is exact under that null and not a discounted figure. The null
itself is narrower than the arithmetic it stands in for: an
i.i.d. sign ensemble on this object measured $0.6455$ of the
spread of a random \emph{multiplicative} one --- a model in the
tradition of \cite{Wintner} --- at $512$ draws
(\texttt{audit\_cn\_coin\_deep.py}).
That measurement rests on one band, one class family and one
statistic, and its own registration says so: it is the second
moment $E[G^2]$ on the classes cut by $3,5,7$ over
$(2^{20},2^{21}]$, not the first-order $m_c-\overline m$ on the
depths cut by $3,5,7,11,13$ over the octave read here. No map
between the two is established anywhere, so the direction is
measured and the size is not: an ensemble that shares $\mu$'s
multiplicativity is the wider of the two, and how much of the
$-11$ that would take back is not fixed by anything in this
paper. The conclusions below survive the direction;
the word \emph{calibrated} would not, and is not used.
Dividing instead by
the permuted floor \eqref{eq:permfloor}, which at depth~$5$ is
$1.1403\cdot10^{-1}$, gives $z\approx-42$. The neighbouring question
--- the numerator against the permutation null \eqref{eq:permnum} ---
has the exact answer $42.42$ (Note~\ref{note:numnull}), and the two
sit close because \eqref{eq:permfloor} is \eqref{eq:permnum} averaged
over the sign null. Reading them as one number quoted conservatively
and boldly would be reading the count-based bar as the truth, and by
Lemma~\ref{lem:permfloor} that bar is what \eqref{eq:permnum} is: at
depth~$5$ it is $1.1356\cdot10^{-1}$, which is
$\mathrm{sd}(Z)/\sqrt{n_c}$ to five figures. The ten draws averaged
$1.1518\cdot10^{-1}$, and that number is a sample where the two
displayed are not. The sign null is the one this paper's statements
are quoted against, throughout. The result files cited here name
it \texttt{lem:coin}, after the statement of the generation they
were written in; that name means this sign null and no further
statement, and it is the null the width comparison earlier in this
note is about.
\end{note}

\begin{lemma}[the floor under the placebo, in closed form]%
\label{lem:permfloor}
Let $c$ be a cell of size $n_c$ in a band of size $n$, let $\pi$ be the
uniformly random permutation of Lemma~\ref{lem:placebo}, and let
$c^\pi$ be $c$ under the labelling $\ell\circ\pi$ --- a uniformly
random $n_c$-subset of the band. Write $\mathrm{Var}_\varepsilon$ for
the variance under the sign null of Lemma~\ref{lem:cellmom} and
$E_\pi$, $\mathrm{Var}_\pi$ for expectation and variance over $\pi$
with the field $Z$ held fixed. Then
\begin{equation}\label{eq:permfloor}
  E_\pi\,\mathrm{Var}_\varepsilon\bigl(m_{c^\pi}-\overline m\bigr)
  \;=\;\frac{n-n_c}{n_c\,(n-1)}\Bigl(1-\frac{Q_{aa}}{n^2}\Bigr),
\end{equation}
and, for the field as observed,
\begin{equation}\label{eq:permnum}
  \mathrm{Var}_\pi\bigl(m_{c^\pi}-\overline m\bigr)
  \;=\;\frac{n-n_c}{n_c\,(n-1)}\,s^2,
  \qquad
  s^2=\frac1n\sum_N\bigl(Z(N)-\overline m\bigr)^2,
\end{equation}
with $E_\varepsilon s^2=1-Q_{aa}/n^2$. Neither right-hand side sees
the cell beyond its size: $Q_{aa}$ is the one term of
\eqref{eq:cellmom} that does not depend on $c$.
\end{lemma}

\begin{proof}
Put $g(N)=V(N)^{-1/2}$ and
$K(N,N')=g(N)g(N')\sum_v\mu^2(v)\Lambda(N-v)\Lambda(N'-v)$, so that
$K$ is the covariance of $Z$ under independent signs and
$Q_{cd}=\sum_{N\in c}\sum_{N'\in d}K(N,N')$. The diagonal is
\[
  K(N,N)=\frac{\sum_v\mu^2(v)\Lambda(N-v)^2}{V(N)}=1,
\]
which is the normalisation of $Z$. It is worth saying why the equality
survives the null. $V(N)=\sum_v\mu^2(v)\Lambda(N-v)^2$ is a function of
the support of $\mu^2$ and of $\Lambda$ alone; the null replaces the
signs $\mu(v)$ on that support and leaves the support itself fixed, so
$V(N)$ is deterministic and unchanged, and the numerator above is
literally $V(N)$. Every use of $K(N,N)=1$ below --- and the whole of
this lemma rests on that diagonal --- is therefore exact rather than
approximate. With
$J_c(N)=\mathbf 1_c(N)-n_c/n$ one has
$m_c-\overline m=n_c^{-1}\sum_NJ_c(N)Z(N)$, so
\begin{equation}\label{eq:JKJ}
  \mathrm{Var}_\varepsilon\bigl(m_c-\overline m\bigr)
  =\frac1{n_c^2}\sum_{N,N'}J_c(N)J_c(N')K(N,N'),
\end{equation}
which is \eqref{eq:cellmom} with its three terms collected. Under
$\pi$ the indicator of $c^\pi$ is that of a uniform $n_c$-subset, so
with $p=n_c/n$, and since $E_\pi J_{c^\pi}=0$,
\[
  E_\pi\bigl[J_{c^\pi}(N)J_{c^\pi}(N')\bigr]
  =p(1-p)\Bigl(\delta_{NN'}-\frac{1-\delta_{NN'}}{n-1}\Bigr).
\]
Taking $E_\pi$ of \eqref{eq:JKJ} and using $\sum_NK(N,N)=n$ and
$\sum_{N\neq N'}K(N,N')=Q_{aa}-n$,
\[
  E_\pi\,\mathrm{Var}_\varepsilon
  =\frac{p(1-p)}{n_c^2}\Bigl(n-\frac{Q_{aa}-n}{n-1}\Bigr)
  =\frac{p(1-p)}{n_c^2}\cdot\frac{n^2-Q_{aa}}{n-1},
\]
which is \eqref{eq:permfloor}. For \eqref{eq:permnum} repeat the
computation with $K$ replaced by the fixed matrix
$(Z(N)-\overline m)(Z(N')-\overline m)$: its diagonal sums to $ns^2$
and its off-diagonal part to $-ns^2$, since $\sum_N(Z-\overline m)=0$.
Finally
$E_\varepsilon\sum_N(Z-\overline m)^2=n-Q_{aa}/n$. The two are one
identity: by Fubini, \eqref{eq:permfloor} is \eqref{eq:permnum}
averaged over the sign null --- the without-replacement variance of a
sample mean, with the population variance replaced by its sign-null
value.
\end{proof}

\begin{note}[what the collapse is]\label{note:collapse}
Lemma~\ref{lem:permfloor} makes the collapse of
Measurement~\ref{meas:placebo} a computation. For a random cell the
floor is, in average over the permutation of the variance, the
without-replacement
formula and nothing else: it decays like $n_c^{-1/2}$, which is exactly the count-based reflex this paper
began by rejecting, with $\mathrm{sd}(Z)$ replaced by
$(1-Q_{aa}/n^2)^{1/2}$ and a finite-population factor attached. That
factor is $\sqrt{(n-n_c)/(n-1)}$ --- the standard finite-population
correction for a sample mean drawn without replacement
\cite{Cochran} --- and it is near $1$ only for a cell
small against its band: at this octave, where $n=10^6$, the two
shallow cells hold $383617$ and $421978$ of it and the factor reads
$0.785$ and $0.760$, against $0.99997$ at depth~$5$. At the
octave $(2\cdot10^6,4\cdot10^6]$, $n=1000000$ and
$1-Q_{aa}/n^2=0.871282$ against $s^2=0.864029$, a ratio of
$0.991676$. \textbf{So the count-based scale is the Rademacher-null
scale for a random cell and is not for the arithmetic cells cut here}
--- the thesis of this paper is a statement about which cells, not
about counts. The
two halves are not equally supported: the first is
Lemma~\ref{lem:permfloor} and holds for every uniformly random subset,
the second is Measurement~\ref{meas:placebo} and covers the depth cells
of one band. A cell cut some third way --- by an interval of the band,
or by the size of $N$ --- is neither, and nothing here says the bar
holds there.

The exact floor of the arithmetic cell exceeds \eqref{eq:permfloor} by
\[
\begin{array}{r|cccccc}
 \text{depth} & 0 & 1 & 2 & 3 & 4 & 5\\\hline
 \mathrm{se}_c\big/\sqrt{\eqref{eq:permfloor}}
   & 104.80 & 42.68 & 70.28 & 44.37 & 16.97 & 3.83
\end{array}
\]
--- the factors quoted above, now read off a formula rather than off
ten draws. Those ten draws stand to $\sqrt{\eqref{eq:permfloor}}$ in
the ratios $0.9968$, $1.0195$, $0.9991$, $1.0175$, $0.9992$,
$1.0101$, and their variances to \eqref{eq:permfloor} itself in
$0.9960$, $1.0420$, $0.9992$, $1.0379$, $0.9999$, $1.0216$; the
largest departure is $1.12$ of the draws' own standard error
(\texttt{lab\_mask\_placebo.py}, L5). Ten draws and a formula agree,
and it is the formula that is quoted.
\end{note}

\begin{note}[the numerator's null is exact]\label{note:numnull}
The second identity supplies what ten draws could not. Under
$\ell\circ\pi$ the numerator $m_{c^\pi}-\overline m$ is the mean of a
uniform $n_c$-subset of the fixed multiset $\{Z(N)\}$ minus its mean:
mean $0$, variance \eqref{eq:permnum} exactly, no draw involved. At
the same octave the true labelling's numerators sit at
\[
\begin{array}{r|cccccc}
 \text{depth} & 0 & 1 & 2 & 3 & 4 & 5\\\hline
 m_c-\overline m & +0.01325 & +0.05191 & -0.03432 & -0.58002
   & -1.99506 & -4.81664\\
 \bigl|m_c-\overline m\bigr|\big/\sqrt{\eqref{eq:permnum}}
   & 11.25 & 47.71 & 16.33 & 106.89 & 102.22 & 42.42
\end{array}
\]
standard deviations of the permutation null. Chebyshev's inequality,
which needs nothing beyond \eqref{eq:permnum}, bounds the permutation
$p$-value of the depth-$5$ numerator by $5.56\cdot10^{-4}$ and the
union over the six cells by $1.28\cdot10^{-2}$. That replaces the
$p\le1/11$ of Measurement~\ref{meas:placebo} by a computation
\emph{for the numerator}. It does not do so for the ratio $z_c$: the
permutation law of the denominator is not given here, and
Measurement~\ref{meas:placebo}'s rank test remains the only statement
about $\max_c|z_c|$ under permutation.
\end{note}

% =====================================================================
\section{What predicts the floor's size}\label{sec:size}

Write $\SS_2(h)$ for the Hardy--Littlewood singular series of the
shift $h$ --- the local density of pairs $(p,p+h)$ --- and, for a cell
$c$ in the band, put
\[
  E_{\mathrm{same},c}[\SS_2]
  \;=\;\frac{1}{n_c(n_c-1)}\sum_{\substack{N,N'\in c\\ N\neq N'}}\SS_2(N-N'),
  \qquad
  E_{\mathrm{all}}[\SS_2]
  \;=\;\frac{1}{n(n-1)}\sum_{\substack{N,N'\in a\\ N\neq N'}}\SS_2(N-N').
\]
Both sums run over \emph{ordered} pairs with the diagonal removed. The
ordering is immaterial, since $\SS_2(-h)=\SS_2(h)$; the removal is not,
since $\SS_2(0)$ is not defined, and every count and every denominator
below is the off-diagonal one. Put
\begin{equation}\label{eq:Dc}
  D_c\;:=\;E_{\mathrm{same},c}[\SS_2]-E_{\mathrm{all}}[\SS_2].
\end{equation}

\begin{measurement}[the floor tracks $D_c$, including its shape]%
\label{meas:Dc}
At the octave $(2\cdot10^6,\,4\cdot10^6]$:
\[
\begin{array}{r|cccccc}
 \text{depth} & 0 & 1 & 2 & 3 & 4 & 5\\\hline
 D_c & 0.302138 & 0.046766 & 0.412504 & 1.052129 & 1.946109
   & 3.171298\\
 \mathrm{se}_c & 0.12400 & 0.04662 & 0.14834 & 0.24178 & 0.33253
   & 0.43685
\end{array}
\]
with correlation $0.9805$ across depths. Neither row is monotone: both
dip at depth $1$, and that is the mechanism made visible. Depth $0$ is
the cell on which \emph{none} of $3,5,7,11,13$ divides $N$, and
excluding a class concentrates the shift just as fixing one does ---
if $3\nmid N$ and $3\nmid N'$ then both lie in $\{1,2\}\bmod3$ and
$3\mid h$ with probability $\tfrac12$ against $\tfrac13$ for a random
pair. Both ends of the depth index concentrate $h$; depth $1$ mixes
the patterns and washes out. Depth is therefore a coarse index for
$D_c$, and the floor follows the same coarse index
(\texttt{lab\_cell\_singular.py}). The control is that run's own, and
is an enumeration rather than a sample: over all $720$ permutations of
the depth labels exactly one reaches $|r|=0.9805$, against a null
median of $0.3475$ and a $95$th percentile of $0.8084$. The
label permutation of Lemma~\ref{lem:placebo} does not supply it, since
that one permutes the labels of $z_c$ and not of this pair.

A correlation over six points cannot tell $\mathrm{se}_c$ from any
other quantity that falls with the depth, and the cell count is one of
those. The ratio settles it and needs no null: $\mathrm{se}_c^2/D_c$
reads $0.05089$, $0.04648$, $0.05334$, $0.05556$, $0.05682$ and
$0.06018$ --- a full width of $25.4\%$ of its mean, which is the
figure the result file prints, and from depth~$1$ upward a monotone
rise of $29\%$ over its smallest value --- while $D_c$ itself moves by
a factor $68$ across the same six cells. That rise is not a property
of the ratio. $\mathrm{se}_c^2$ carries the diagonal of $K$, where
$K(N,N)=1$, and $D_c$ does not: its averages run over pairs $N\ne N'$,
and $\SS_2(0)$ is not defined. The diagonal is
$n_c^{-2}\sum_NJ_c(N)^2=(1-n_c/n)/n_c$ and has no partner in the
denominator. At the deepest cell it is $7.82\%$ of $\mathrm{se}_c^2$,
and what it contributes there, $+0.004706$, exceeds the whole printed
rise from depth~$4$ to depth~$5$, $+0.00336$. That contribution is
depth~$5$'s own share and not the step between the two depths; the
cited run states it as such.

No single factor carries $\mathrm{se}_c^2$ to $D_c$'s convention.
\eqref{eq:JKJ} is the
doubly centred mean $E_{\mathrm{same},c}[K]-2E_{c,a}[K]+E_{\mathrm{all}}[K]$,
which is three sums over three index sets with three different pair
counts --- $n_c^2$, $n_cn$, $n^2$. Each term has to lose its own
diagonal and be divided by its own count, and $K(N,N)=1$ makes that
exact. Done that way the row reads $0.05089$, $0.04646$, $0.05334$,
$0.05553$, $0.05665$, $0.05704$: a full width of $19.8\%$ rather than
$25.4\%$, monotone from depth~$1$ still, and a step from depth~$4$ to
depth~$5$ of $+0.00038$. \textbf{The rise survives the match; what
shrinks is its size} --- nine times smaller at the last step, and a
rise of $23\%$ from depth~$1$ rather than $29\%$. Removing only the
diagonal gives $-0.00112$ there and scaling that by $n_c/(n_c-1)$
gives $-0.00031$; both are over-corrections and neither is the
comparison to make. The two quantities
are one functional of two kernels: \eqref{eq:JKJ} is the doubly
centred mean $n_c^{-2}\sum J_cJ_cK$, and $D_c$ is that same mean with
$K$ replaced by the singular series of the shift. The centring is not
cosmetic. For $K$ the cross term is not the band mean --- it runs
$0.124636$ to $0.140049$ against $E_{\mathrm{all}}=0.128718$, and
dropping it would misstate the variance by factors $0.469$ to $2.120$
across these six depths (\texttt{lab\_mask\_placebo.py}). For the
singular series it is: $E_{c,a}$ departs from $E_{\mathrm{all}}$ by at
most $0.004192$, so the two-term difference \eqref{eq:Dc} and the
doubly centred form agree to within one per cent at five of the six
depths --- the largest of those five being $0.59\%$ at depth~$3$ ---
and to within five per cent at depth~$1$, where $D_c$ is
smallest (\texttt{lab\_cell\_singular.py}). Both bands quoted here have
the same numerator, the three-term variance $\mathrm{se}_c^2$, and
differ only in the denominator: against $D_c$ as \eqref{eq:Dc} defines
it the band is the $25.4\%$ above, and against the doubly centred
$B_c$ it is $29.7\%$.
A near-constant ratio through the origin is what would follow if $K$
and the singular series differed across the band by a factor that the
double centring leaves alone. That is not proved here: the ratio is
measured, not derived. What the measurement does settle is the
alternative it was taken against --- a quantity that merely decreases
with the depth produces no such ratio. What it does not settle
is the other direction. The band is read across depths at one
octave; across scale the same ratio falls over the three octaves on
which $D_c$ is measured. Divided by the sampled $D_c$ it reads
$7.9\%$, $0.9\%$, $9.4\%$, $9.2\%$, $9.6\%$, $15.2\%$, and the
$0.9\%$ at depth~$1$ was read here as that depth being the exception.
It is not: divided by the $D_c$ this measurement computes over all
pairs, the six read $9.5\%$, $9.6\%$, $9.5\%$, $9.6\%$, $9.8\%$,
$15.3\%$. \textbf{The exception was the estimator's} --- at depth~$1$
the pair average carries more noise per octave than the drift it was
said to show. What remains is a drift of about half the matched width
at five of the six cells and rather more at the deepest. Both inputs are printed here
(\texttt{lab\_cell\_singular.py}, which forms the
quotient from its own $D_c$ and the variance row of
\texttt{lab\_cell\_floor.txt}). So $D_c$ fixes the
ratio's depth profile at a fixed scale and not its scale
dependence, and the floor itself is not scale-invariant
(Measurement~\ref{meas:exponent} fits it a nonzero
exponent) while
$D_c$ is: $D_c$ is therefore not all of what fixes the
floor's size.
\end{measurement}

The proof of the proposition below sums a weight along an arithmetic
progression. What it needs of the weight is only bounded variation, and
that is worth isolating, because it is the one step where the shape of
the weight --- as opposed to its arithmetic --- does any work.

\begin{lemma}[summing along a progression against an envelope]%
\label{lem:envtv}
Let $w:\mathbb Z\to\mathbb R$ be supported on $|h|<B$ with total
variation $\mathrm{TV}(w)=\sum_h|w(h+1)-w(h)|$, and let $q\ge1$,
$a\in\mathbb Z$. Then
\[
  \sum_{h\equiv a\ (q)} w(h)
  \;=\;\frac1q\sum_h w(h)\;+\;O\bigl(\mathrm{TV}(w)\bigr),
\]
with an absolute implied constant --- in particular uniform in $q$, in
$a$, and in $B$.
\end{lemma}

\begin{proof}
Partition all of $\mathbb Z$ --- not the support --- into the blocks
$I_j$ of $q$ consecutive integers determined by $a$, and let $h_j$ be
the representative of the class in $I_j$. Every block is then complete,
and all but finitely many meet the support not at all and contribute
zero to both sides. So
\[
  \sum_{h\equiv a\,(q)}w(h)-\frac1q\sum_h w(h)
  \;=\;\sum_j\Bigl[w(h_j)-\frac1q\sum_{h\in I_j}w(h)\Bigr]
  \;=\;\sum_j\frac1q\sum_{h\in I_j}\bigl[w(h_j)-w(h)\bigr],
\]
and $|w(h_j)-w(h)|\le V_j$ for every $h\in I_j$, where $V_j$ is the
variation of $w$ within $I_j$. The inner average is therefore at most
$V_j$, and the $V_j$ are sums of $|w(h+1)-w(h)|$ over disjoint sets of
$h$, so $\sum_j V_j\le\mathrm{TV}(w)$. The implied constant is $1$.
\end{proof}

Nothing in it is arithmetic: $q$ is arbitrary, and the bound does not
improve with $q$ nor degrade with it. That is exactly why it survives
the primes comparable to the band length, which are equidistributed
modulo nothing.

\begin{proposition}[the size mechanism, asymptotically scale-free]%
\label{prop:scaleinv}
We establish an asymptotic statement and not an identity. Let
$Q=3\cdot5\cdot7\cdot11\cdot13$, let $a$ be the even integers of an
interval of length $B$, and put $\ell=B/2Q$, so that each residue
class modulo $2Q$ meets $a$ in $\ell+O(1)$ terms. Let $R_c$ be a
nonempty set of $m$ residue classes modulo $2Q$ and
$c=\{N\in a:\ N\bmod 2Q\in R_c\}$ the cell they cut. Every depth cell
is of this form --- depth fixes how many of the five primes divide
$N$ and not the residues themselves, so $R_c$ is the set of classes
of the $\binom{5}{d}$ divisibility patterns of depth $d$ --- and $a$
itself is the case in which $R_c$ is all $Q$ even classes. Write
\[
  A(h)=\prod_{2<p\mid h,\ p\le13}\Bigl(1+\frac1{p-2}\Bigr),
  \qquad
  C_2=\prod_{p>2}\frac{p(p-2)}{(p-1)^2},
  \qquad
  \kappa=2C_2\prod_{p>13}\Bigl(1+\frac1{p(p-2)}\Bigr),
\]
so that $A$ is a function of $h\bmod Q$, and put
\[
  \mathcal A_c=\frac1{m^2}\sum_{r,r'\in R_c}A(r-r'),
  \qquad
  \mathcal A_a=\frac1{Q^2}\sum_{r,r'}A(r-r'),
\]
the second sum over all pairs of even classes modulo $2Q$. These are
rational numbers determined by $R_c$ alone: they depend neither on
$B$ nor on where the interval sits. Then, with $D_c$ as in
\eqref{eq:Dc} and for every fixed $\eta\in(0,1)$,
\begin{equation}\label{eq:scaleinv}
  D_c\;=\;\kappa\,(\mathcal A_c-\mathcal A_a)
  \;+\;O\Bigl(\frac1\ell\Bigr)
  \;+\;O_\eta\Bigl(\frac{B^{\eta}}{\ell}\Bigr)
  \;+\;O_\eta\bigl(B^{\eta-1}\bigr).
\end{equation}
The implied constant in the first error term is absolute and those in
the other two depend on $\eta$ alone; in particular none of them
depends on the cell $c$, on $m$, on $n_c$, on $B$, or on the position
of the interval, and the parameter set $\{3,5,7,11,13\}$ enters only
through $Q$, fixed throughout. The $\eta$-dependence is not uniform as
$\eta\to0$ --- the constant comes from the Euler products of Rankin's
bound in the proof and blows up there --- so $\eta$ is fixed in
advance and the statement is one about each fixed $\eta$ separately.

Consequently $D_c\to\kappa(\mathcal A_c-\mathcal A_a)$ as
$B\to\infty$ with $R_c$ held fixed. Writing $D_c(B)$ for the value of
\eqref{eq:Dc} in a band of length $B$, define the one-step exponent
\begin{equation}\label{eq:ec}
  e_c(B)\;=\;-\,\frac{\log\bigl|D_c(2B)/D_c(B)\bigr|}{\log 2}.
\end{equation}
Where $\mathcal A_c\neq\mathcal A_a$ the convergence gives
$e_c(B)\to0$; where $\mathcal A_c=\mathcal A_a$ it gives only
$D_c\to0$ and predicts no exponent. No rate is claimed for either
limit, and nothing here asserts equality at a fixed band.
Measurement~\ref{meas:maincoef} writes $\mathcal A_c$ and
$\mathcal A_a$ as $E_{\mathrm{same},c}[A]$ and $E_{\mathrm{all}}[A]$,
and Measurement~\ref{meas:secondcell} writes the first
$E^{\mathrm{ex}}_c[A_S]$; the pair average of $A$ over the cell itself
differs from $\mathcal A_c$ by the first error term.
\end{proposition}

\begin{proof}
For $N,N'$ in a band, $\SS_2(N-N')$ depends on $N-N'$ only through the
local factors at primes dividing $h=N-N'$. Depth does \emph{not} fix
the residue of $N$ modulo each of $3,5,7,11,13$ --- it fixes only how
many of them divide $N$, so a depth-$d$ cell is the union of the
$\binom{5}{d}$ divisibility patterns of that size, and
Measurement~\ref{meas:Dc} records the mixing this causes. What the
argument needs is weaker and does hold: each pattern fixes, at each of
the five primes, the admissible residue set, so within a pattern the
law of $h$ modulo each of them is determined; and, the even $N$ of a
band being equidistributed modulo $3\cdot5\cdot7\cdot11\cdot13$ to
$O(1/\ell)$ per class, the joint residue-class structure within a
pattern is determined by the product modulus. The word is meant in the
Chinese-remainder sense and not the probabilistic one: the admissible
set modulo $3\cdot5\cdot7\cdot11\cdot13$ is the product of the
admissible sets at the five primes, so the joint law modulo the product
is determined and not only the five marginals. No independence of
random variables is asserted or used. The weights with
which the patterns enter a depth cell are their densities among the
even $N$ of the band, which are the same at every scale. The induced
law of $h$ modulo those primes is therefore determined by the cell
alone up to the resolution of the cell itself. That resolution is
$O(1/\ell)$ and not $O(1/n_c)$: a depth cell meets each joint residue
class in about $\ell=B/2Q$ terms, and it is that count, not the cell's
total, that the empirical law of the residues is read from. The same
counting governs the divisor sampling below --- the resolution is
$O(1/\ell)$ for both and not $O(1/n_c)$ --- and no cancellation is
claimed for either.

Write $Q=3\cdot5\cdot7\cdot11\cdot13$ and split
$\SS_2(h)=2C_2\,A(h)\,T(h)$, where
$A(h)=\prod_{2<p\mid h,\ p\le13}(1+\frac1{p-2})$ is a function of
$h\bmod Q$ and $T(h)=\prod_{p\mid h,\ p>13}(1+\frac1{p-2})$ carries
everything above $13$. What has to be shown is that $T$ contributes
the same factor to both averages of \eqref{eq:Dc}, whatever the cell.
That is the hardest step of the proof and not a hypothesis readable
off the construction. It is not enough that each prime above $13$ is
separately equidistributed: $\SS_2$ reads them all at once, and the
primes comparable to the band length are equidistributed modulo
nothing.

Expand $T$ over divisors. Writing $g(d)=\prod_{p\mid d}\frac1{p-2}$
and letting $d$ run over the squarefree integers all of whose prime
factors exceed $13$, $T(h)=\sum_{d\mid h}g(d)$. Every such $d$ is
coprime to $2Q$; and the shifts that occur are $h=N-N'$ with both $N$
even, so they lie in classes modulo $2Q$ and not modulo $Q$. For a
residue $a$ modulo $2Q$ the Chinese remainder theorem gives
\[
  \#\{h:0<|h|<B,\ h\equiv a\ (2Q),\ d\mid h\}
  =\frac{B}{Qd}+O(1),
\]
and $d>B$ divides no $h$ of the band at all, since $h=N-N'$ with both
$N$ inside a band of length $B$.

The weight with which $h$ enters a pair average is not a function of
$|h|$ alone. For $E_{\mathrm{same},c}$ it counts the $N\in c$ with
$N-h\in c$, and membership of $c$ is fixed by $N\bmod Q$, so the count
reads $h\bmod Q$ as well: if $3\mid h$ the shift preserves divisibility
by $3$ and otherwise it does not. What is true, and what the argument
uses, is that the weight factors --- a function of $h\bmod Q$, constant
on each residue class, times an envelope in $|h|$. That envelope has
total variation $O(1/B)$, and Lemma~\ref{lem:envtv} applied to it gives
$P(d\mid h)=1/d+O(1/B)$ uniformly in $d\le B$. The mixture over
classes is what the pattern weights supply. Every cell's pair law is such a
mixture, whence
\[
  E[\SS_2]=2C_2\,E[A]\prod_{p>13}\Bigl(1+\frac1{p(p-2)}\Bigr)
  \;+\;\mathcal E,
  \qquad
  \mathcal E=O(1/\ell)+O_\eta\bigl(B^{\eta}/\ell\bigr)
    +O_\eta\bigl(B^{\eta-1}\bigr)\ \ (\eta>0).
\]

All three parts of $\mathcal E$ have to be said. None of them depends
on the cell: $\ell=B/2Q$ is a band parameter, the same at depth~$0$ and
at depth~$5$. A cell-dependent bound would have to be an $O(1/n_c)$,
and the counting below does not support one.
Which part is the larger is settled in $B$ --- the divisor-sampling
term exceeds the first by $B^{\eta}$ and the third by $2Q$ --- and is
unsettled only in the implied constants, which the conclusion does not
need. The $O(1/B)$ above is the error of the
envelope, not of the cell. A cell is not a random sample: it is a
union of residue classes modulo $2Q$, and every $d$ here is coprime to
$2Q$, so the count that served the envelope serves the cell as well ---
taken with the cell's own parameters. A depth cell is $m$ classes, each
meeting the band in $\ell=B/2Q$ terms, with $m\ell=n_c$; only the
deepest cell of this field is a single class. Over one pair of classes
the pairs with $d\mid h$ number $\ell^2/d+O(\ell)$, so over the $m^2$
pairs of classes the error is $O(m^2\ell)=O(n_c^2/\ell)$ and the
proportion departs from $1/d$ by $O(1/\ell)=O(Q/B)$ --- uniformly in
$d$, with no gain from the size of $d$. Summed against $g$ that is
$\ell^{-1}\sum_{d\le B}g(d)$, which the Rankin bound below puts at
$O_\eta(B^{\eta}/\ell)$; $\ell$ is a fixed multiple of $B$, so this is
$O_\eta(B^{\eta-1})$ again.

The counting gives $O(1/\ell)$ and not $O(1/n_c)$, and the two differ
by the factor $m$, which is large at every cell but the deepest. The
stronger form would need the $m^2$ class-wise errors to cancel rather
than accumulate. Nothing here establishes that, and it is not claimed. This is the
same $O(1/\ell)$ that the equidistribution above carries, arrived at by
the same count of a cell against its residue classes; the two are one
resolution seen twice and not two errors. Both are bounds rather than
errors, and which of them dominates is settled in $B$ and unsettled
only in the implied constants. Neither varies by cell. The conclusion
needs no answer to that, only that both go to zero, and $\ell$ grows
with the band.

The last part is the two truncations, which Rankin's trick bounds.
For $s>0$, $\sum_{d\le B}g(d)\le B^{s}\prod_{p>13}(1+\frac1{(p-2)p^{s}})$,
the product converging because $\sum_p p^{-1-s}$ does; this is
$O_\eta(B^{\eta})$ at $s=\eta$, and reaches $O_\eta(B^{\eta-1})$ only
after the per-$d$ error $O(1/B)$ multiplies it. For $t<1$,
$\sum_{d>B}g(d)/d\le B^{-t}\prod_{p>13}(1+\frac1{(p-2)p^{1-t}})$, the
product converging because $\sum_p p^{-2+t}$ does; at $t=1-\eta$ this
is $O_\eta(B^{\eta-1})$ directly. The two reach the same order by
different routes, and the two exponents cannot be collapsed into one:
with $s=t$ the bounds read $B^{s-1}$ and $B^{-s}$, which meet only at
$s=\tfrac12$, and $B^{-1/2}$ is far weaker than $B^{\eta-1}$. Nor does
$\eta$ pay to be taken small without limit: $\log$ of the product is
of order $\log(1\!/\!\eta)$, so the bound is best near
$\eta=1/\!\log B$ and worsens on either side of it. The two exponents
stay two, but at $s=\eta$ and $t=1-\eta$ the two \emph{products}
coincide: both read $\prod_{p>13}(1+\frac1{(p-2)p^{\eta}})$, and the
two convergence criteria both reduce to $\sum_p p^{-1-\eta}$. There is
one $\eta$-constant in this proof, not two.

Every constant here is uniform in the cell, which is what lets
$\mathcal E$ be one bound for a whole band. $A$ is exactly constant on
each pair of residue classes modulo $2Q$ --- $h$ is fixed there --- and
is bounded by $\prod_{2<p\le13}(1+\frac1{p-2})=128/33$, a number with
no cell in it. On one class pair the ordered pairs $(i,j)$ in
$[0,\ell)^2$ with $i-j$ in a fixed class modulo $d$ number
$\ell^2/d+O(\ell)$ with an absolute constant, since the difference
count is a tent of total variation $2\ell$. And the $m^2$ class pairs
divide out identically against $n_c^2=m^2\ell^2$: the same relative
error appears at $m=5760$ and at $m=1$. Nothing in the bound names the
cell's pattern set.

What is \emph{not} uniform is $\mathcal E/D_c$, and it is that ratio
the exponent needs. At the octave Measurement~\ref{meas:scaleinv}
reads, $D_c$ runs from $0.046766$ at depth~$1$ to $3.171298$ at
depth~$5$, a factor of $67.8$; one and the same $\mathcal E$ therefore
buys $67.8$ times less at the first than at the last, at every band.
So the proposition is not equally informative at the six cells, and
depth~$1$ is the cell where it says least --- the same depth whose main
term the proof cannot certify is nonzero.

The factor the primes above $13$ contribute is therefore one and the
same constant $\kappa/2C_2$ in both averages, and it leaves the
difference of the two pair averages of $A$. Each of those is the
average of a function of $h\bmod Q$ over the pairs of a union of
residue classes, and the classes meet the band in $\ell+O(1)$ terms
each, so each pair average is $\mathcal A_c$, respectively
$\mathcal A_a$, up to $O(1/\ell)$ with an absolute constant, $A$ being
bounded by $128/33$. Hence
\[
  D_c=\kappa\,(\mathcal A_c-\mathcal A_a)+\mathcal E,
\]
which is \eqref{eq:scaleinv}. The main term is fixed by $R_c$ and is
the same at every scale, so \eqref{eq:Dc} converges to it. Convergence
is additive and an exponent is multiplicative: if the limit is nonzero
then $D_c(2B)/D_c(B)\to1$ and the one-step exponent \eqref{eq:ec}
tends to $0$, while if $\mathcal A_c=\mathcal A_a$ the proposition
says only that $D_c\to0$ and predicts no exponent at all. No derivative
is taken anywhere: $D_c$ is defined band by band, and what the
convergence controls is the ratio of consecutive bands, which is what
\eqref{eq:ec} reads. The difference $\mathcal A_c-\mathcal A_a$ is a
finite sum of rationals over the residues modulo $2Q$ and the pattern
weights, so it is decidable cell by cell without any measurement.
Measurement~\ref{meas:maincoef} decides it: the difference is a
nonzero rational at all six depths, so the exponent this proposition
predicts is defined at all six.
\end{proof}

\begin{measurement}[the main term, computed]\label{meas:maincoef}
$A(h)$ depends only on which of $3,5,7,11,13$ divide $h$, a cell is a
union of residue classes modulo $2Q$, and $p\mid h$ for $h$ between two
classes exactly when the two agree modulo $p$. So
$E_{\mathrm{same},c}[A]-E_{\mathrm{all}}[A]$ is a finite double sum
over pattern pairs weighted by $m(S)=\prod_{p\notin S}(p-1)$, and it is
computed in exact rational arithmetic
(\texttt{lab\_maincoef.py}): $635653709/2854051200$,
$10743253/313945632$, $318004704/1037775739$, $296100736/378194817$,
$13823072/9546537$ and $78976/33033$ at depths $0$ to $5$. \textbf{None
is zero}, which is the hypothesis the proof needs. At depth~$5$ the
cell is a single class, so $Q$ divides every shift in it and
$E_{\mathrm{same},c}[A]$ is exactly
$\prod_{2<p\le13}(1+\frac1{p-2})=128/33$.

Multiplied by $2C_2\prod_{p>13}(1+\frac1{p(p-2)})$ --- a constant with
no fitted parameter in it, and one that is not irrational: the same
term-by-term cancellation that gives $E_\infty[\SS_2]=2$ below reduces
it to
\[
  2C_2\prod_{p>13}\Bigl(1+\frac1{p(p-2)}\Bigr)
  \;=\;2\prod_{p\in\{3,5,7,11,13\}}\frac{p(p-2)}{(p-1)^2}
  \;=\;\frac{11011}{8192}\;=\;1.3441162109375
\]
exactly --- these predict $D_c$ at the octave
$(2\cdot10^6,4\cdot10^6]$:
\[
\begin{array}{r|cccccc}
 \text{depth} & 0 & 1 & 2 & 3 & 4 & 5\\\hline
 \text{predicted} & 0.299361 & 0.045996 & 0.411876 & 1.052351
   & 1.946236 & 3.213542\\
 \text{measured} & 0.302138 & 0.046766 & 0.412504 & 1.052129
   & 1.946109 & 3.171298\\
 \text{residual}/\mathrm{se} & 2.0 & 0.6 & 0.5 & 0.2 & 0.1 & 260.8
\end{array}
\]
Six numbers are predicted from none, across cells spanning a factor of
$68$. \emph{Four of the six sit inside two sampling errors of the
measured value}, depth~$0$ sits at $2.0$, and the sixth is depth~$5$ at
$260.8$. Both of the last two are accounted for below and neither by a
fitted quantity. Over all $720$ relabellings of the depth
labels the count of cells agreeing within two sampling errors reaches
the observed four twice, $p=2/720$; the six predictions span a factor
of $70$, so a wrong pairing is refused by size alone, and what the null
shows is that the agreement is carried by which cell gets which
number.

\textbf{Depth~$5$'s miss is not a residual but a computable quantity,
and computing it is what the cell's structure allows.} That cell is a
single residue class modulo $2Q$, hence an arithmetic progression
segment $N_0+2Qi$; every same-cell shift is $h=2Qk$ with $0<|k|<n_c$;
and every $d$ in the divisor expansion is coprime to $2Q$, so $d\mid h$
exactly when $d\mid k$. \emph{No $d\ge n_c$ divides any same-cell shift
at all}, while the Euler factor sums over every $d$. The resulting
deficit is a finite sum, $\Delta=8.505279\cdot10^{-3}$ at $n_c=67$, and
$2C_2\cdot\tfrac{128}{33}\cdot\Delta=0.043558$ against an observed
residual of $0.042244$ --- agreement to $3.1$ per cent, from a
computation with no free parameter and no asymptotics. It takes
depth~$5$ from $260.8$ sampling errors to $8.1$.

Calling that miss ``the size of $1/\ell$'' would instead be the reading
the discussion below withdraws, and unfalsifiable besides: with an
unnamed implied constant no discrepancy could refuse it.

\textbf{It is also the same quantity Measurement~\ref{meas:scaleinv}
records as depth~$5$'s spread.} $\Delta$ falls with $n_c$ --- it is
$1.297902\cdot10^{-2}$, $8.505279\cdot10^{-3}$ and
$5.208268\cdot10^{-3}$ at $n_c=34,67,134$ --- so it induces a drift in
$D_5$ across the three octaves of $1.24$ per cent, against the $1.28$
that measurement finds. \emph{The two are not on one denominator}: the
first is taken over the uncorrected prediction, the second over a mean
of measured values, and the cited run prints both readings rather than
one. What the comparison supports is the order of the effect, not an
agreement to the second digit. \emph{So the largest of the six spreads is not
a failure of scale invariance: it is the proposition's own error term,
computed.} Both figures were predicted before this computation existed
by an independent reading of the construction, which put the residual
at $1.35$ per cent and claimed only the order for the spread.

The constant cuts the singular series at $13$, and an off-by-one in
which primes are absorbed into $C_2$ would move it by $1/143$, or
$0.70$ per cent. It is not there: such an error is a uniform
multiplicative factor, and one of that size would show at $11.5$
sampling errors at depth~$4$ alone.

\textbf{The band average has a limiting value in closed form, and the
run estimated it instead.} Write $E_{\mathrm{all},B}$ for the pair
average over the band of length $B$ --- the quantity
\eqref{eq:Dc} subtracts --- and $E_\infty$ for its limit. The average
of the singular series over shifts is classical \cite{Gallagher}, and
the error term in it has been sharpened since \cite{Pintz}; what is
needed here is only the leading constant. That constant is $1$ when the
average is taken over \emph{all} shifts, and the average needed here is
twice it. The reason is that $\SS_2(h)=0$ for odd $h$ --- the local
factor at $2$ vanishes, since $n$ and $n+h$ cannot both be odd primes
above $2$ --- while every shift occurring here is even, $h=N-N'$ with
$N$ and $N'$ both even. Averaging a function that vanishes on half the
integers over the half where it does not doubles its mean. Taking
$P(p\mid h)=1/p$ for every odd $p$ gives
$E_\infty[\SS_2]=2C_2\prod_{p>2}(p-1)^2/(p(p-2))$, and
$C_2=\prod_{p>2}p(p-2)/(p-1)^2$ is the reciprocal of that product term
by term, so
\[
  E_\infty[\SS_2]=2 .
\]
\emph{This is a limiting value and not the band's own average}, which is
$1.999986$ at the band Measurement~\ref{meas:secondcell} reads; the two
differ by $1.4\cdot10^{-5}$, and that difference is the band's
equidistribution error, discussed there.

The run's $E_{\mathrm{all},B}$, recoverable from its own table as
$E_{\mathrm{same},c}-D_c$ and the same number in all six rows, is
$1.998755$: low against the limit by $1.2\cdot10^{-3}$, which is a
hundred times the finite-band difference and is therefore the
\emph{sampling} error of that average. It is a single additive offset
shared by every $D_c$, since each is $E_{\mathrm{same},c}$ minus that
one number; it is one sampling error at five cells and eight at the
sixth, because depth~$5$'s $4422$ pairs are enumerable where $10^{12}$
are not. It is not a truncated prime product: the sieve here runs over
every prime to $8\cdot10^6$, and a truncation would in any case leave
depth~$5$ exact, every same-cell shift there being free of prime
factors above $61$.

Substituting the limiting value for the sampled one --- writing
$D'_c:=E_{\mathrm{same},c}[\SS_2]-2$, which is \eqref{eq:Dc} only up to
the band's own $1.4\cdot10^{-5}$ --- the prediction reads
$\tfrac{11011}{8192}E_{\mathrm{same},c}[A]-2$, one estimated input
fewer, and
the row closes:
\[
\begin{array}{r|cccccc}
 \text{depth} & 0 & 1 & 2 & 3 & 4 & 5\\\hline
 D'_c=E_{\mathrm{same},c}-2 & 0.300893 & 0.045521 & 0.411259 & 1.050884
   & 1.944864 & 3.170053\\
 \text{predicted} & 0.299361 & 0.045996 & 0.411876 & 1.052351
   & 1.946236 & 3.169984\\
 \text{residual}/\mathrm{se} & -1.13 & +0.36 & +0.45 & +1.07 & +1.16
   & -0.43
\end{array}
\]
\textbf{All six sit inside $1.2$ sampling errors, and nothing is
fitted}: the $2$ is a theorem, the depth-$5$ entry carries the finite
divisor deficit above, and the constant is the singular series' own.
Depth~$5$ has gone from $260.8$ to $-0.43$.

The computation is post hoc as to timing: it was written after the
$D_c$ row was measured. It is not a fit --- the difference is forced by
the residue structure and the constant by the definition of the
singular series --- but a registered rule of it was refuted. The rule
predicted the worst relative deviation at depth~$5$; it is at depth~$1$
($1.65$ per cent). \textbf{The rule was written on the wrong
statistic}: the measured $D_c$ is a sampled estimator whose sampling
error this repository already prints beside it, and against that error
depth~$1$ is $0.6$ and depth~$5$ is $260.8$. Read on the right
quantity the claim holds; read as registered it does not, and it is
recorded refuted.
\end{measurement}

\begin{measurement}[a second cell shape, predicted first]\label{meas:secondcell}
Everything above is a statement about one partition. The computation
of Measurement~\ref{meas:maincoef} runs for any set of odd primes, so a
second partition can be \emph{predicted before it is measured}, which
the first can no longer be given.

Cells indexed by depth over $S=(3,5,7)$, so $Q'=105$ and
$\ell=B/2Q'$ is $9524$ at this octave against the first partition's
$66.6$. The constant is $K=2\prod_{p\in S}p(p-2)/(p-1)^2=175/128$,
exactly rational, and the predictions
$K\,E_{\mathrm{same},c}[A_S]-2$ were computed and committed
(\texttt{lab\_secondcell\_predict.py}) before any measurement of this
partition existed. The measurement (\texttt{lab\_secondcell2.py},
registered in \texttt{lab\_secondcell2\_predict.py}) is exact: a cell
is a set of
positions in the band, so the pair average of $\SS_2$ over it is one
autocorrelation of the cell indicator, with no sampling, no seed and no
standard error. The pair counts reproduce $n_c(n_c-1)$ exactly at all
four cells.
\[
\begin{array}{r|cccc}
 \text{depth} & 0 & 1 & 2 & 3\\\hline
 m & 48 & 44 & 12 & 1\\
 \text{main term, uncorrected} & 0.295736 & 0.126485 & 0.825521 & 2.375000\\
 \text{fully corrected} & 0.295286 & 0.126068 & 0.824967 & 2.374142\\
 \text{predicted, }1/m\text{ corrected} & 0.295726 & 0.126476 & 0.825475 & 2.374142\\
 \text{exact} & 0.295708 & 0.126454 & 0.825436 & 2.374142\\
 \text{difference} & -0.000018 & -0.000022 & -0.000039 & -10^{-7}
\end{array}
\]
The first two rows are the endpoints T1 gates on and are printed so
that the rule can be checked from the table rather than from the file:
the deficit is negative, so the interval at each depth is
$[\text{fully corrected},\,\text{uncorrected}]$. The row the third line
predicts is neither endpoint --- it is the $1/m$ diluted correction,
which is what T2 scores.
Of the three registered rules two hold and one is void. The exact
value lies between the main term and the fully corrected prediction at
all three of the cells where the two differ, which is the rule that
gates, and the fitted exponent stays inside the spread the third rule
allowed. \textbf{So the computation transfers to a cell shape it was
not built on.} Depth~$3$ is a
consistency check and not a score: there the prediction and the
measurement are the same finite sum, as the registration recorded
before the run.

\emph{One} rule was registered against the evidence then available.
The sampled thirty-draw mean at depth~$2$ was $0.825796$, above T1's
upper bound of $0.825521$ by about $1.2$ standard errors of that mean,
and the sealed file records both that lean and that it was not
decisive. The exact value is inside the interval. So that is one pass
against a stated lean --- not two, and not against a prediction of
failure: the sealed file says ``leans, not decisive'' of the first and
``may fail'' of the second, and the second is void rather than passed.

\textbf{The correction the prediction carries should not be modelled at
all, and this measurement is what shows it.} The prediction diluted the
divisor deficit by $m$, the number of residue classes a cell unions.
Fitting an exponent instead, $(\text{undiluted})/m^{\alpha}$, gives
$0.715$, $0.689$ and $0.756$ at $m=48,\,44,\,12$. The registration
fixed in advance that a value strictly between $0$ and $1$ says the
model is wrong rather than that the exponent is that value, so
\textbf{no exponent is claimed}: a dilution exists and $1/m$ is not
it.

It does not need to. The deficit is the same autocorrelation with
$T=\prod_{p\mid h,\,p>7}(p-1)/(p-2)$ in place of $\SS_2$, so it is
computed rather than fitted --- and computing it exposes a fourth term
that no fit could have separated:
\[
  D'_c=\underbrace{K\,E^{\mathrm{ex}}_{c}[A_S]-2}_{\text{main term}}
        +\underbrace{K\bigl(E_c[A_S]-E^{\mathrm{ex}}_{c}[A_S]\bigr)}
          _{\text{pattern weight}}
        -\underbrace{2C_2E_c[A_S]\Delta_c}_{\text{divisor deficit}}
        +\underbrace{2C_2\mathrm{Cov}_c(A_S,T)}_{\text{covariance}},
\]
where $E^{\mathrm{ex}}_c$ is the exact rational the prediction is built
from and $E_c$ the measured pair average.
\[
\begin{array}{r|cccc}
 \text{depth} & 0 & 1 & 2 & 3\\\hline
 \text{main term} & 0.2957357 & 0.1264850 & 0.8255208 & 2.3750000\\
 \text{pattern weight} & -4.55\cdot10^{-6} & -6.75\cdot10^{-6}
   & -3.99\cdot10^{-6} & 0\\
 \text{deficit} & -1.32\cdot10^{-5} & -1.23\cdot10^{-5}
   & -5.35\cdot10^{-5} & -8.58\cdot10^{-4}\\
 \text{covariance} & -1.02\cdot10^{-5} & -1.17\cdot10^{-5}
   & -2.69\cdot10^{-5} & 0\\
 \text{sum} & 0.2957078 & 0.1264543 & 0.8254364 & 2.3741421
\end{array}
\]
and the sum reproduces the measured value at every cell. No exponent
and no free parameter survive anywhere in it.

The pattern-weight term is the departure of the empirical class weights
from the idealised ones --- which is
Proposition~\ref{prop:scaleinv}'s \emph{other} error term, the
equidistribution one, measured here for the first time --- and the
covariance is the failure of $E[A_ST]$ to factor; \textbf{both are
exactly zero at a single-class cell, which is the only kind either
partition ever calibrated the deficit on}, so a fitted exponent was
carrying two quantities that cannot appear where it was calibrated.

\textbf{No power law in $m$ fits these three}, which is a firmer
reason to drop the exponent than the absence of a mechanism. The
exponent implied between $m=48$ and $m=44$ is $0.04$ against $0.91$
between $m=44$ and $m=12$; the deficits at the first pair are in the
ratio $1.004$ where a power fitted to the second requires $1.083$. The
measurement is exact, so that is not noise. A monotone dependence on
$m$ is not excluded --- three points cannot exclude one --- but no
single power reaches them.

What the second cell shape measures, once the exponent is dropped, is
\textbf{how tight Proposition~\ref{prop:scaleinv}'s own error term is}.
At depth~$0$ the three measured terms are $5.68$, $4.36$ and $1.95$ per
cent of the \emph{number} $1/\ell=1.05\cdot10^{-4}$, and complete
cancellation of the class-wise errors --- what the diagonal class pairs
alone would give, which the proof declines to claim --- would put the
deficit at $4.03$. The proposition's bound is that number times an
unnamed constant and a divisor sum above three, as the discussion after
its proof records, so these are \emph{upper} bounds on the fraction of
the bound the measurement realises: the bound is at least three times
looser than the figures suggest, and the conclusion that it is loose
survives a fortiori.

\textbf{The measured deficit exceeds the diagonal-only figure} at every
multi-class cell, by $41$, $30$ and $16$ per cent at $m=48,\,44,\,12$,
so the off-diagonal pairs add deficit rather than cancelling it. Those
three are quoted and not summarised: an eight per cent change in $m$
between the first two moves the excess by thirty-eight, so the excess
is no more a function of $m$ than the deficit is.

One caution belongs with those numbers, and it turns into a check. The
registration fixed $D'_c=E_{\mathrm{same},c}[\SS_2]-2$, and the rules
are scored on that convention; but $2$ is the \emph{limiting} value
$E_\infty$, while \eqref{eq:Dc} subtracts the band's own average
$E_{\mathrm{all},B}$, which is $1.999986$ here. The two differ by
$1.39\cdot10^{-5}$, \emph{larger than two of the four terms above}, so
$D'_c$ and $D_c$ are two quantities under one name at this band size.
Substituting the band's own average removes that same
$1.39\cdot10^{-5}$ at every cell --- half the departure from the main
term at depth~$0$ and under two per cent of it at depth~$3$, since a
constant offset cannot take a fixed share off departures spanning a
factor of thirty. \textbf{T1 holds under both conventions}, so no
verdict turns on the choice; the margin at depth~$0$ is
$1.40\cdot10^{-5}$ against a shift of $1.39\cdot10^{-5}$, which is thin
enough that another systematic of the same size would break it.

\textbf{And the band's own departure decomposes into the same three
terms, at a cell where the main term is known in advance to be zero.}
The band is all $105$ classes modulo $2Q'$, so it is a cell of this
partition --- the one with $m=105$ --- and
$K\,E^{\mathrm{ex}}[A_S]-2=(175/128)(256/175)-2=0$ exactly. Its
$-1.38928\cdot10^{-5}$ is then $-2.37\cdot10^{-6}$ of pattern weight,
$-5.83\cdot10^{-6}$ of divisor deficit and $-5.69\cdot10^{-6}$ of
covariance, summing to the measured value to one part in $10^{10}$.
It is the only check here whose answer was fixed before the
computation.

Because no exponent is claimed, \emph{nothing here is to be applied to
the first partition's row}. That row is the main term with no
correction of any kind. A correction computed the way this measurement
computes one would move its depth-$4$ entry, whose cell unions
$34$ classes; the entry is quoted uncorrected and the correction is not
transported, since a law fitted on one partition and applied to another
is exactly the transport this measurement was built to test rather than
to assume.

\textbf{The rule that tested the $1/m$ dilution is recorded neither
held nor refuted, but \emph{void}}, and the reason is worth more than
the rule. Its tolerance was $10^{-4}$; the models it had to separate
differ by $1.8$ to $4.6\cdot10^{-5}$ at these $m$, so both pass it. A
test the alternative also passes carries no evidence in either
direction, and filing it as held would be worse than filing it as
refuted, since a held rule is counted as support and there is none. The
tolerance was set without asking what it had to discriminate; that
question belongs before a rule is registered, not after it is scored.
\end{measurement}

Of the proof's errors it is the cell-sized one that is worth reading
against the measurement, and this is where the sizes belong. The
divisor-sampling term does \emph{not} vary with the depth: it is
$O_\eta(B^{\eta}/\ell)$ with $\ell=B/2Q$, one bound for every cell of
an octave.

\textbf{No percentage can be read off it.} The term carries an implied
constant that the proof does not name, and the per-$d$ error $O(1/\ell)$
enters summed against $g$, so the bound is $\ell^{-1}\sum_{d<B}g(d)$
and that divisor sum is not $O(1)$ --- it grows like $\log B$, and at
$B=2\cdot10^6$ the single and double prime terms alone put it above
three. Reading the term as $1/\ell$, which at that octave is one and a
half per cent, is that bound with both the constant and the divisor
sum set to $1$, and it cannot then be compared with the two per cent
that Measurement~\ref{meas:scaleinv} registers as its cap: the proof
supplies no cap for such a comparison to use.

What survives is the comparison's direction, and it survives a
fortiori. The measurement caps the relative spread of $D_c$ across the
three octaves at two per cent and finds spreads from $0.0001$ at
depth~$0$ to $0.0128$ at depth~$5$; the proof's bound on the error at
each octave is larger than that and unquantified. \textbf{So the
proof's bound is weaker than the measurement at every cell.} It does
not resolve what the exact spreads do, and nothing here should be read
off it about which cell drifts and which does not. In particular it
cannot be read against the sampled spreads to conclude that whatever
moves the two refused depths is not an error term of this proof: the
sampled spreads are the estimator's, the refusals with them, and the
quantity's own spreads are the ones above.

\begin{measurement}[the exponent, measured]\label{meas:scaleinv}
The exponent is discrete, and it is defined by ratios and not by a
derivative: it is the one-step exponent $e_c(B)$ of \eqref{eq:ec},
with $D_c(B)$ the value of \eqref{eq:Dc} computed in the band of
length $B$, the sign being fixed by the convention $|D_c|\sim B^{-e}$,
so that a positive $e$ means a decaying $D_c$. The $e$ tabulated below
is minus the least-squares slope of $\log|D_c|$ against $\log B$ over
the three
octaves. Their abscissae are the octave midpoints, spaced by exactly
$\log 2$, so that quantity is the mean of the two one-step exponents;
nothing is fitted that the ratios do not already determine.

Measured over $(10^6,2\cdot10^6]$, $(2\cdot10^6,4\cdot10^6]$ and
$(4\cdot10^6,8\cdot10^6]$ at $400000$ sampled pairs per cell, the
exponents are
\[
\begin{array}{r|cccccc}
 \text{depth} & 0 & 1 & 2 & 3 & 4 & 5\\\hline
 e & +0.012633 & +0.065546 & +0.000936 & +0.002904 & +0.000241
   & -0.008428
\end{array}
\]
--- zero to three decimals at depth $4$, within $0.003$ at depths $2$
and $3$, and $-0.008$ at depth $5$, against $+0.013$ and $+0.066$ at
depths $0$ and $1$. The two large ones are where $D_c$ is small: the
spread across the three octaves runs from $0.0008$ at depth $2$ to
$0.037$ at depth $5$, the last being the finite-cell drift of a
cell of a few tens of elements, while $D_c$
itself runs from $0.045$ to $3.17$, so the relative spread the exponent
reads is large only where $D_c$ is. At ten times the sample the two
move to $-0.000879$ and $-0.016226$, so depth~$0$ resolves to zero and
depth~$1$ flips sign --- the signature of noise rather than of a trend.
That resampling is post hoc: it was run after M2 and M4 had been
scored, and it does not lift either verdict. At the higher sample
depth~$1$ still reads $e=-0.016226$ against M4's $|e|<0.01$ and a
spread of $0.0390$ against M2's two per cent, so both stay refused
as the code scores them.

The registration's own wording and the code's scoring of it differ, and
this paper reports the stricter --- the code's --- throughout.
Supplement~S1 gives both readings and the verdicts each produces.

Sampling was not necessary. $D_c$ is a deterministic functional of the
cell, and the sum over all ordered pairs off the diagonal is one
autocorrelation per cell:
$\sum_{N\neq N'\in c}\SS_2(N-N')
=\sum_{h\neq0}\SS_2(h)\,(\mathbf 1_c\star\mathbf 1_c)(h)$, the
diagonal being excluded on both sides because $\SS_2(0)$ is undefined;
what is dropped is the term $(\mathbf 1_c\star\mathbf 1_c)(0)=n_c$.
Computed that way
the six read
\[
\begin{array}{r|cccccc}
 \text{depth} & 0 & 1 & 2 & 3 & 4 & 5\\\hline
 \text{spread} & 0.0001 & 0.0004 & 0.0002 & 0.0002 & 0.0022 & 0.0128\\
 e & -0.000059 & -0.000285 & -0.000136 & -0.000164 & -0.001584
   & -0.009202
\end{array}
\]
--- \textbf{six of six inside M2's two per cent and six of six inside
M4's $|e|<0.01$}, the shallow four by three orders. The two depths the
registered rules refuse are the two where the sampling error is the
largest fraction of $D_c$, and the refusals are a property of that
estimator and not of $D_c$. The verdicts stand as scored, on the
estimator the run registered. The deepest cell keeps the largest
spread, $1.28$ per cent, which Measurement~\ref{meas:maincoef}
computes as $1.24$ on its own denominator --- the divisor deficit of a
cell of a few tens of elements, and the only one of the six that the
estimator was reading rather than adding to. Supplement~S1 records the
registered rules, what each was scored on, and where in the result
files each refusal can be reproduced.

The sampling error of the pair average runs $0.001183$ to $0.001370$
at depths $0$ to $4$ --- $0.001358$ at depth~$0$, and lowest at
depth~$4$, where the cell is smallest --- and $0.000162$ at
depth~$5$, so at depth~$1$, where
$D_c=0.046766$ is six times smaller than any other, that error is the
largest fraction of $D_c$ anywhere in the table. That is what the
refusals at depth~$1$ rest on, and the exact computation above says so
in the other direction: remove the estimator and the spread there is
$0.0004$. Depth~$0$ is refused by the exponent rule alone and not by
this: there the sampling error is $0.45\%$ of $D_c$, and what fails is
$e=+0.012633$ against $|e|<0.01$, which the tenfold resampling
resolves to $-0.000879$ and the exact computation to $-0.000059$.
Depth~$1$ is quiet under
the sign null ($z=+1.11$ in Measurement~\ref{meas:placebo}; in
absolute value the third smallest of the six, behind depth~$0$'s
$+0.1069$ and depth~$2$'s $-0.2314$), but that is a
statement about the sign null: under the permutation null of
Note~\ref{note:numnull}, which this paper calls the stronger
instrument, the same cell reads $47.71$ --- third of the six, behind
the $106.89$ and $102.22$ at depths~$3$ and~$4$, and above the $42.42$
at depth~$5$ that the paper calls the signal. The cell is therefore not one that
carries no effect, and the two refusals are not excused by saying
it does (\texttt{lab\_cell\_singular.py},
\texttt{lab\_mask\_placebo.py}).
\end{measurement}

Proposition~\ref{prop:scaleinv} settles one thing and deliberately not
another. Whatever produces the decay of the effect with $N$, it is not
the arithmetic excess that produces the floor's size, since that
excess does not decay at all. The decay exponent of the effect itself
is not determined here. At the three shallowest depths no exponent is
fitted, and the amplitude there is not uniformly below the floor
across the range: at depth~$1$ it falls monotonically from $+2.71$ at
the smallest octave to $+0.69$ at the largest, and two octaves cannot
settle it unless they are the two smallest. The floor test asks of each scale separately whether
the amplitude is present; whether it is falling is a question about
the scales together, and failing the first does not answer the
second.

% =====================================================================
\section{Summary}

\begin{itemize}
\item Lemma~\ref{lem:cellmom}: the fluctuation of any cell mean under
the sign null is exactly
$Q_{cc}/n_c^2-2Q_{ca}/(n_cn)+Q_{aa}/n^2$; all three terms are needed
and all three are one convolution each.

\item Observation~\ref{obs:coh} and Note~\ref{note:notanerrorbar}:
that quantity decays like $(\log N)^{-1/2}$ at the shallow cells,
where Note~\ref{note:threeterms}'s fraction is constant, and not like
$n_c^{-1/2}$ anywhere. The exponent itself is measured at all six
depths and the deepest fits $0.107306$, which the form does not
predict. A
count-based interval is narrower than the exact Rademacher-null scale
by a factor $9.3$ across a factor $128$ in $N$, and absolutely by
factors of $7.3$ to $158$ at the top octave, the largest of them at the
shallowest cell.

\item Lemma~\ref{lem:placebo} and Measurement~\ref{meas:placebo}:
neither the cell effect nor the floor is reproduced under label
permutation, so both read the correspondence and not the cell sizes
--- the floor least expectedly, since one expects a floor to be a
property of the sizes. The floor
is a second measurement of the same correspondence. What it collapses
to is the count-based bar (Lemma~\ref{lem:permfloor}), so the two
floors are two questions and not a bracket; $z_c$ is quoted against
the sign null throughout.

\item Proposition~\ref{prop:scaleinv}: what is established is an
asymptotic statement, not an identity. $D_c$ equals a cell-determined
constant up to $O(1/\ell)$ and $O_\eta(B^{\eta-1})$, both of which
vanish with the band and neither of which carries a named implied
constant; so no rate and no percentage may be read from it, and its
consequence --- that this quantity is not what makes the effect decay
--- holds in the limit and not at any fixed band. The argument is
given in \S\ref{sec:size}. It is the one labelled statement here
that has had neither a machine check nor a reading by a subject
expert, and its hardest step is named as such in its own proof.
Two of the six depths are refused by the registered rules, and both
refusals belong to the sampling estimator rather than to $D_c$: the
same measurement computes $D_c$ exactly, by one autocorrelation per
cell, and done that way all six depths sit inside both registered
caps, the shallow four by three orders. The verdicts stand as scored,
on the estimator the run registered; Supplement~S1 gives them. It is the proposition's own error term, not a failure of it.

\item Measurement~\ref{meas:maincoef}: that quantity's main term is a
finite rational over the residue classes the cells are unions of, it
is nonzero at every depth --- which is the hypothesis
Proposition~\ref{prop:scaleinv} needs and does not otherwise have ---
and multiplied by a constant with no fitted parameter it reproduces
the measured profile, four of six inside two sampling errors across
cells spanning a factor of $68$. \textbf{It is the only quantity here
predicted from the arithmetic alone}: Lemma~\ref{lem:cellmom} also
gives a size in closed form, but from the field's own values, and this
takes no input from the field at all. It is post hoc in timing. A registered rule of it was refuted: it named the worst
relative deviation at the wrong cell, having been written on the
unstandardised statistic while the standardisation was printed beside
it.

\item The mechanism --- coherence of $u_c(v)$ across the cell,
\emph{together with} a cell correlated with the arithmetic that
coherence comes from --- is generic. Sharing the arithmetic input is
not by itself enough: Lemma~\ref{lem:permfloor} gives a random cell of
the same band a floor whose permutation average is exactly the
count-based one, and that cell shares the input.
\end{itemize}

All preregistered checks, including the failed ones and the one
recorded void, are reported with their verdicts in Supplement~S1
(Appendix~\ref{sec:repro}). Several numerical patterns conjectured in
the course of the work were rejected there; none is used as an
assumption in anything stated above.

% =====================================================================
\begin{thebibliography}{99}

\bibitem{HL} J.-J.~Huang and H.~Li, \emph{On the connection between
the Goldbach conjecture and the Elliott--Halberstam conjecture},
arXiv:2005.03811 [math.NT]; v1 (May 2020), v2 (August 2022). Also
published in a Springer volume,
\texttt{doi:10.1007/978-3-030-67996-5\_17}. References here are to the
arXiv version, which we have read, and not to the published version,
which we have not.

\bibitem{Mir49} L.~Mirsky, \emph{The number of representations of an
integer as the sum of a prime and a $k$-free integer}, Amer.\ Math.\
Monthly \textbf{56} (1949), 17--19. Theorem~1, taken at $k=2$: the
count is $\prod_{p\nmid n}(1-p^{1-k}/(p-1))\operatorname{Li}n
+O(n/\log^Hn)$ for every $H>0$, whose local factor at $k=2$ is the
$1-1/(p(p-1))$ used above.

\bibitem{Gallagher} P.~X.~Gallagher, \emph{On the distribution of
primes in short intervals}, Mathematika \textbf{23} (1976), 4--9;
corrigendum, Mathematika \textbf{28} (1981), 86. What is used here is
the theorem on the average of the singular series over $k$-tuples. We
have not obtained the original; the statement is used in the form
in which \cite{Pintz} states and reproves it.

\bibitem{Pintz} J.~Pintz, \emph{On the singular series in the prime
$k$-tuple conjecture}, arXiv:1004.1084v1 [math.NT], 2010. What is used
above is Theorem~1, $\SS_{\mathcal H}(H)=1+O(\varepsilon)$ for
$H\ge\exp(k^{1/\varepsilon})$, applied with $\mathcal H$ a single
element: there $\SS(\mathcal H)=1$ and
$\SS(\mathcal H\cup\{h\})$ is the singular series of the shift, so the
average over shifts tends to $1$. Equation~(1.5),
$\sum_{|H|=k,\,H\subset[1,H]}\SS(H)\sim H^{k}$, is Gallagher's global
average and is \emph{not} what is used: its $k=1$ case is the sum over
singletons, and a singleton has $\SS=1$ identically, so that case is
an identity carrying no information about shifts. The paper itself
draws the distinction, noting that $\SS_{\mathcal H}(H)\to1$ ``yields
more information about the singular series than the global average''.
\S1 also records $\SS(\mathcal H\cup\{h\})=0$ for odd $h$, which is
the vanishing that makes the average over even shifts twice the
average over all.

\bibitem{Cochran} W.~G.~Cochran, \emph{Sampling Techniques}, 3rd ed.,
Wiley, 1977. (Finite-population correction for the variance of a
sample mean without replacement.)

\bibitem{Kish} L.~Kish, \emph{Survey Sampling}, Wiley, 1965. (Design
effect and effective sample size for correlated observations.)

\bibitem{LR} E.~L.~Lehmann and J.~P.~Romano, \emph{Testing Statistical
Hypotheses}, 3rd ed., Springer, 2005. (Permutation and randomisation
tests conditional on the observed configuration.)

\bibitem{Wintner} A.~Wintner, \emph{Random factorizations and
Riemann's hypothesis}, Duke Math.\ J.\ \textbf{11} (1944), 267--275,
\texttt{doi:10.1215/S0012-7094-44-01122-1}.

\end{thebibliography}

% =====================================================================
\section*{Disclosure statement}

No potential competing interest, financial or non-financial, was
reported by the author. No funding was received for this work. An AI
assistant was used in preparing it; what tool, how, and why is set out
in Appendix~\ref{sec:repro}, together with what was checked by machine,
what by computation, and what by neither.

\section*{Data availability statement}

The code behind every figure printed here, the output it produced, and
the Lean~4 development are openly available at
\texttt{https://github.com/soke91/arithmetic-convolution-fields}, in
the directory \texttt{P4-coherent-cell-floor/}. The figures quoted are
those of the release named in Appendix~\ref{sec:repro}. No data derived
from human or animal subjects were used; every quantity here is
computed from the integers.

% =====================================================================
\appendix
\section{The scale of the normaliser}\label{sec:scale}

Observation~\ref{obs:coh}'s heuristic reads the scale $N\log N$ off the
asymptotic size of $V(N)$. That asymptotic is stated and proved here so
that this paper depends on no unpublished manuscript; nothing else in
the paper uses it, every computation dividing by $V(N)$ as computed.

\begin{proposition}[the scale of $V$]\label{prop:V}
Let $V(N)=\sum_{v<N}\mu^2(v)\Lambda(N-v)^2$ and
$W(N)=\sum_{w<N}\Lambda(w)^2$, and put
$\AAA(N)=\prod_{q\nmid N}\bigl(1-1/(q(q-1))\bigr)$. Then
$V(N)=W(N)\AAA(N)(1+o(1))$, and hence $V(N)\sim\AAA(N)\,N\log N$, the
second form using $W(N)\sim N\log N$.
\end{proposition}

\begin{proof}
$\Lambda(w)^2$ is supported on prime powers with weight $(\log p)^2$,
so with $w=N-v$,
\[
  V(N)=\sum_{p^k<N}(\log p)^2\mu^2(N-p^k)
      =\sum_{p<N}(\log p)^2\mu^2(N-p)+O(\sqrt N\log^2N),
\]
a $(\log p)^2$-weighted count of squarefree shifted primes. Its local
density at a prime $q$ is $1$ when $q\mid N$ --- there $q^2\mid N-p$
forces $q\mid p$, hence $p=q$, a single term --- and
$1-1/\varphi(q^2)=1-1/(q(q-1))$ when $q\nmid N$, the bad class
$p\equiv N\ (\mathrm{mod}\ q^2)$ being a unit class. The count is then
Mirsky's theorem on squarefree values of shifted primes
\cite{Mir49}, a statement at the endpoint $x=N$. The weight needs no
uniform-in-$x$ partial summation, but the cut has to move: for
\emph{fixed} $\varepsilon$ the ratio $(\log p)^2/(\log N)^2$ on
$p>N^{1-\varepsilon}$ lies in $((1-\varepsilon)^2,1]$ and does not
tend to $1$, so that choice gives $(1+O(\varepsilon))$ and not
$(1+o(1))$. Take instead $\varepsilon=\varepsilon(N)\to0$ slowly
enough that $\varepsilon\log N/\log\log N\to\infty$ ---
$\varepsilon=1/\log\log N$ serves. Then
$(1-\varepsilon)^2=1+o(1)$, so $(\log p)^2=(1+o(1))(\log N)^2$ on
$p>N^{1-\varepsilon}$, while
$\pi(N^{1-\varepsilon})=o(N/\log N)$ still
absorbs the smaller primes whatever their weight, so the weight is
$(1+o(1))(\log N)^2$ on all but $o(N/\log N)$ of the terms. The second
form follows from $W(N)\sim N\log N$, which is partial summation from
$\theta(x)\sim x$ and is the only analytic input either form takes.
\end{proof}

\noindent
The local factor here is $\AAA$ and not the Hardy--Littlewood singular
series $\SS$: both depend on $N$ only through the primes dividing it,
but $\SS$ through a product \emph{over} those primes and $\AAA$ through
a product over the primes \emph{not} among them, so a substitution of
one for the other is easy to make and hard to see.

\section{Reproducibility}\label{sec:repro}

The code, the output it produced, and the Lean development are at

\begin{center}
\texttt{https://github.com/soke91/arithmetic-convolution-fields}
\end{center}

\noindent in the directory \texttt{P4-coherent-cell-floor/}. That
address is mutable. The code and output every figure here is read from
are those of release \texttt{v1.0.1}, commit \texttt{b2f365d}, and a
number checked against the branch head may have been recomputed since.
The pin is on \texttt{code/} and \texttt{results/} and not on this
document, which cannot cite a commit that contains it. Each entry
in the table below is run as
\texttt{python code/<script>.py} from that directory, and each writes
the file of the same name in \texttt{results/}. The scripts import only
the standard library and NumPy; they never import one another. Five of
them do read a result file that another script produces, so the packet
carries those producers too and the order matters:
\texttt{lab\_cell\_singular} reads \texttt{lab\_cell\_floor} and
\texttt{lab\_mask\_placebo}; \texttt{lab\_secondcell2\_predict} reads
\texttt{lab\_secondcell} and \texttt{lab\_secondcell\_predict};
\texttt{lab\_secondcell2} reads both \texttt{predict} files;
\texttt{lab\_maincoef} reads \texttt{lab\_cell\_singular} and
\texttt{lab\_secondcell2}; and \texttt{audit\_cn\_coin\_deep} reads
\texttt{audit\_cn\_class} and \texttt{audit\_cn\_multnull\_deep}. Run in
that order, every input exists before the script that reads it.

The figures printed here were produced under Python~3.12.10 and
NumPy~2.2.5. \texttt{lab\_maincoef.py} and the two \texttt{predict}
files compute in exact rational arithmetic and their output is
reproducible bit for bit on any platform; the rest compute in double
precision, where a different BLAS may move the last printed digit of a
summed quantity. No figure quoted in this paper is stated to more
digits than that distinction supports.

Four of the runs exit with a nonzero status. That is their registered
behaviour and not a failure: a run whose own preregistered rule was
refuted says so in its exit code as well as in its output, and the
refutations are reported in Supplement~S1.

The Lean development is in \texttt{P4-coherent-cell-floor/lean/}, with
Mathlib pinned in \texttt{lake-manifest.json} and the toolchain in
\texttt{lean-toolchain}; \texttt{lake build} reproduces
\texttt{axioms.txt}, which is the axiom report quoted above.

\begin{center}
\begin{tabular}{lll}
\hline
Script & Result file & Supports\\
\hline
\texttt{lab\_cellmom\_montecarlo.py} &
 \texttt{lab\_cellmom\_montecarlo.txt} & Measurement~\ref{meas:mc}\\
\texttt{lab\_cell\_floor.py} & \texttt{lab\_cell\_floor.txt}
 & Measurement~\ref{meas:exponent}\\
\texttt{lab\_mask\_placebo.py} & \texttt{lab\_mask\_placebo.txt}
 & Measurement~\ref{meas:placebo}\\
\texttt{lab\_cell\_singular.py} & \texttt{lab\_cell\_singular.txt}
 & Measurements~\ref{meas:scaleinv},~\ref{meas:Dc}\\
\texttt{lab\_maincoef.py} & \texttt{lab\_maincoef.txt}
 & Measurement~\ref{meas:maincoef}\\
\texttt{lab\_secondcell2.py} & \texttt{lab\_secondcell2.txt}
 & Measurement~\ref{meas:secondcell}\\
\texttt{audit\_cellfloor\_sdz.py} & \texttt{audit\_cellfloor\_sdz.txt}
 & Note~\ref{note:notanerrorbar}\\
\texttt{audit\_cn\_coin\_deep.py} & \texttt{audit\_cn\_coin\_deep.txt}
 & Note~\ref{note:floorsignal}'s $0.6455$\\
\hline
\end{tabular}
\end{center}

\subsection*{Use of an AI assistant, and what checked what}

This work was done in close collaboration with an AI assistant, which
carried much of the derivation and all of the computation. The author
takes full responsibility for every claim made here.

\emph{Which tool, and how.} The manuscript in its present form was
prepared with Anthropic's Claude, model Opus~5, run through the Claude
Code command-line interface. It was used to derive and check the
algebraic identities, to write and run every script in the accompanying
code, to write the Lean development, and to draft and revise this text.
Earlier stages of the work, before the material reported here took its
present shape, used earlier models of the same family; the repository
does not record which, and this paper does not rely on anything from
those stages that has not since been rederived and rerun. The reason
for use is throughput on a body of computation and case analysis larger
than the author would otherwise have carried out, and the cost of that
choice is the one this note is about: what a reader can check does not
depend on trusting the tool.

What follows says what was checked and by what, rather than asserting
that it was checked.

\emph{Machine-verified.} Lemma~\ref{lem:cellmom}'s algebraic identity,
Lemma~\ref{lem:permfloor}'s closed form and Lemma~\ref{lem:placebo}
are proved in Lean~4 against Mathlib; the axiom report is in the code
packet and each depends on \texttt{propext}, \texttt{Classical.choice}
and \texttt{Quot.sound} and on nothing else. \textbf{Proposition~\ref{prop:scaleinv}
is not machine-verified and has had no reading by a subject expert.}
Its statement is written out in Lean and left unproved, so the axiom
report carries \texttt{sorryAx} on that line and the absence is
checkable rather than a matter of trust. The finite half of its
argument is verified separately; the analytic half, the Euler product
and its convergence, is not.

\emph{Checked by computation.} Every figure quoted is produced by a
named run in the accompanying code, and a check in the build refuses a
figure that no cited run prints. Every registered rule is reported with
its verdict, including the thirteen the runs did not meet and the one
recorded void. The predictions of Measurement~\ref{meas:secondcell} were
committed to the repository before the measurement that scores them
existed.

\subsection*{Supplement S1: what rests on what}

Lemma~\ref{lem:contrast}, Lemma~\ref{lem:cellmom},
Lemma~\ref{lem:placebo}, Lemma~\ref{lem:permfloor},
Lemma~\ref{lem:envtv} and Proposition~\ref{prop:scaleinv} are proved
and use no computation; those six are every proof in the paper.
Lemmas~\ref{lem:contrast} and~\ref{lem:envtv} use no arithmetic
either: the first is a quadratic form in an arbitrary covariance, the
second a summation bound for an arbitrary modulus, and neither is
among the statements formalised in Lean.
The scripts were written against the manuscript of the generation they
belong to and their headings name its statements, not this paper's:
\texttt{prop:coh}, \texttt{prop:placebo}, \texttt{rem:floorsignal} and
\texttt{rem:sdzrange} are Observation~\ref{obs:coh},
Lemma~\ref{lem:placebo}, Note~\ref{note:floorsignal} and
Note~\ref{note:notanerrorbar} under earlier names, and
\texttt{lem:coin} is the sign null, as Note~\ref{note:numnull}
records. \texttt{rem:copiedfloor}, \texttt{rem:levelmeas},
\texttt{rem:mcratios} and \texttt{rem:cncoindeep} are four further
statements of that manuscript which this paper does not carry; the
runs are cited here for their numbers and nothing above rests on
those four.
The Supports column names one statement per run and there are more.
\texttt{lab\_cell\_floor.py} also carries
Measurement~\ref{meas:notzero}, the paragraph of
Measurement~\ref{meas:mc} that reads
Note~\ref{note:threeterms}'s structural constant at the top octave
--- the one that measurement itself ends by attributing to this file
--- the table of
Observation~\ref{obs:coh}, Notes~\ref{note:threeterms}
and~\ref{note:notanerrorbar}, and two items of
\S\ref{sec:notclaimed} --- the one on the shallow depths not being
uniformly quiet, and the one on Observation~\ref{obs:coh}'s looseness
factors, whose two inputs are that file's $Q_{cc}/n_c^2$ row and its
heuristic $0.049540$.
\texttt{lab\_mask\_placebo.py} also carries
Notes~\ref{note:floorsignal}, ~\ref{note:collapse}
and~\ref{note:numnull} and the \(\mathrm{se}_c\) row of
Measurement~\ref{meas:Dc}. The width ratio $0.6455$ inside
Note~\ref{note:floorsignal} is not that file's: it comes from
\texttt{audit\_cn\_coin\_deep.py},
the sixth run of the table above. Counting a named statement as
resting on a run when it prints a figure that only that run produces,
five rest on the first file --- Observation~\ref{obs:coh},
Note~\ref{note:threeterms} and Measurements~\ref{meas:mc},
\ref{meas:exponent} and~\ref{meas:notzero} --- and six on the second:
Measurement~\ref{meas:placebo}, Notes~\ref{note:floorsignal},
\ref{note:collapse} and~\ref{note:numnull}, the $\mathrm{se}_c$ row of
Measurement~\ref{meas:Dc}, and the sign-null amplitudes that
Measurement~\ref{meas:scaleinv} reads off it. The unit is stated because
the count turns on it: counting the items of \S\ref{sec:notclaimed} and
the paragraphs that carry no figure of their own alongside the named
statements would make the first eight. Under a unit that can be checked
the two numbers are five and six.
The measurements state the band, the cell index and the statistic they
were taken of, and each was run with its decision rule fixed in
advance. Two qualifications belong beside that sentence.
The first is that not every registered rule can fail. N3 of
Measurement~\ref{meas:mc} reads the Monte-Carlo precision at the
smaller draw count against the figure an earlier printing quoted; the
precision is a function of the draw count alone, and the draw count is
a constant of the script, so N3 evaluates the same arithmetic on every
run and cannot be refuted by any input. It is an identity, and it is
counted among the registered rules below.
The second is that one registration block carries text written after
its run. The heading of \texttt{lab\_mask\_placebo.py} reads ``written
before this script was run'', and L5 inside it argues for its own cap
by citing the factors $3.8$ to $105$ that the run produced, while L4
records in the same block what the run failed to test. The rules
themselves were fixed first --- what came later is the reasoning
printed beside them --- but a reader who takes the heading literally
is reading two layers as one.
One quantity a measurement leans on carries no rule of its
own: the spread of $\mathrm{se}_c^2/D_c$ that
Measurement~\ref{meas:Dc} calls settling --- $25.4\%$ as printed,
$19.8\%$ under the matched convention. The $19.1\%$ that the cited
file also prints is the diagonal-out row, which is an over-correction
and which that measurement withdraws. What is registered there is
M3, the correlation; the ratio came out of the same run and is gated
by nothing. The enumeration over all $720$ relabellings that scores M3
\emph{is} registered --- that file declares it in its heading --- and
the ten-fold resampling of Measurement~\ref{meas:scaleinv} is post hoc
and says so. The permutation control of Lemma~\ref{lem:placebo} was
\emph{not} run against each cell-indexed claim. It was run once: ten
label permutations on the single octave $(2\cdot10^6,\,4\cdot10^6]$,
with the cells indexed by depth. Of the six cell-indexed measurements
the rest carry no control by \emph{that} permutation ---
Measurement~\ref{meas:Dc} carries its own, an enumeration over all
$720$ orderings of the depth labels --- and two of the cited files say
so where they say it best --- one records that the placebo ``is not
run here'', and the other that the control the placebo file supplies
is not the control that statement needs, since it permutes labels for
a different statistic. The control is not claimed for the others, and
what the evidence files say is what stands.
\subsection*{Supplement S1: the registered rules and their verdicts}

Six of the ten runs record a rule their own output did not meet. The
unit of the count is the pair (run, label), because the labels are
local to their files and repeat across them: \texttt{N1} names one
rule in \texttt{lab\_cellmom\_montecarlo} and a different one in
\texttt{lab\_maincoef}, and \texttt{F1} and \texttt{F2} likewise name
different rules in \texttt{lab\_secondcell\_predict} and
\texttt{audit\_cn\_coin\_deep}. Under that unit thirty-eight rules
carry a verdict --- twenty-four hold, thirteen refuted, one void ---
and one further label, W3, sits under the same registration heading as
W1, W2 and W4 and is reported without a verdict, so thirty-nine labels
are registered in all. The thirteen are named here because a
reader who does not open the files would otherwise take the runs for
clean, and because in every case what the rule asked for is not what
the statement above says. The seventh,
\texttt{lab\_cellmom\_montecarlo}, met all four of its rules, and
two of those four score a $60$-draw quotation that an earlier printing
of this work made and this paper does not: the table above is at
$2000$ draws, and the precision Measurement~\ref{meas:mc} reasons from
is computed in the text rather than gated by any rule of that file.
The ratios that measurement states are checked in a block the file
marks post hoc.

\texttt{lab\_cell\_floor} registered that the fitted exponent $b$ in
$\mathrm{se}\sim N^{-b}$ over the eight octaves, and the quantity
$1/(2\langle\log N\rangle)$ beside it, would match the values published
earlier (C4); that the row of $Q_{cc}/n_c^2$ would have minimum
$0.124$ and maximum $0.307$ at the top octave (C3); and that the four
ratios of C1 would behave the same way at the octave
$(2\cdot10^6,4\cdot10^6]$ (C2). C4 is refuted by one line of four: the
three fitted $b$ agree with what was published, and
$1/(2\langle\log N\rangle)$ comes out $0.036038$ against a published
$0.0358$ --- the value used in this paper is the new one, and nothing
above turns on the third decimal. C3 is refuted on its minimum, which
measures $0.121208$ against the registered $0.124$: the row's minimum
sits at depth~1 rather than at depth~0, so the published figure was
reading the wrong end. Two of C3's three clauses fail, not one: the
maximum holds at $0.307074$ against the registered $0.307$, but the
heuristic reads $0.049540$ against $0.049$, which is $0.000540$ away
and so outside the $0.0005$ the rule allowed --- it fails as
registered, by a margin smaller than the last digit of the published
value it was registered against. C2 is
refuted at depth~$5$ alone, which reads $0.557654$ at that octave
against $0.551516$ at the top; that is the sentence in
Measurement~\ref{meas:mc} which says depth~$5$ does not.
\texttt{lab\_cell\_singular}
registered four things and three fail: that $D_c$ is positive and
increasing in the depth (M1); that it is invariant across octaves ---
a spread within two per cent (M2) and a fitted exponent with
$|e|<0.01$ (M4), a threshold its result file does not print; and that
the floor tracks $D_c$ across depths, a correlation above $0.8$ (M3),
which holds at $0.9805$ and is the control
Measurement~\ref{meas:Dc} rests on. Only M1 is about the depth, and its failure is the
sentence in Measurement~\ref{meas:Dc} that neither row is monotone;
M2 and M4 are about scale and fail only where $D_c$ is small, as
recorded above.
\texttt{lab\_mask\_placebo} registered that the floor would survive
permutation within a tenth (L3) --- it does not, the collapse being
the subject of Measurement~\ref{meas:placebo} and the reason
Note~\ref{note:floorsignal} and Lemma~\ref{lem:permfloor} were
written --- and two things at once as L4: that
the true labelling's $z_c$ is monotone decreasing in depth, and that
under permutation the rank correlation between depth and $z_c$ has
$|\mathrm{mean}|<0.3$. L4 is refuted on the first --- $z$ rises by
$1.0065$ from depth~$0$ to depth~$1$ --- while the second holds, the
permuted mean Spearman being $-0.0114$ against a true $-0.9429$. \texttt{audit\_cellfloor\_sdz}
registered that its printed pair would be reproduced under one of two
readings (W1) and that the ratio would grow with the cell (W2); neither
does, and the convention that settles W1 is stated in
Note~\ref{note:notanerrorbar}. \texttt{audit\_cn\_coin\_deep}
registered four things about the two sign ensembles and two fail:
F3 asked the i.i.d. ensemble to cover the true value at the
deepest cell and no draw of it does, and F4 asked the two nulls to
agree in width within a factor $1.5$ and they do not --- the
$0.6455$ that Note~\ref{note:floorsignal} quotes is that
disagreement. F1 and F2 hold. This run is cited here for the width
ratio alone, and its two refutations are the reason that ratio is
worth quoting.

Seven of the thirteen refuted a prediction about a row this paper
prints --- its shape across depths (C3, M1, L4, W2) or its
constancy across octaves (C2, M2, M4) --- and the statements
above report what the rows do. Two more, F3 and F4, are about the
sign null's own reach rather than about this field, and are
named above. A tenth, N3, registered which of six cells carries the
worst deviation and named depth~$5$; its own file records the
refutation together with the reason, that the rule was registered on
an unstandardised quantity and so asked about a different question
from the one it appears to ask. Nothing above rests on it. The other
three
did more than that, and are listed here because the count alone would
hide them. W1 registered that the pair $(5.8,\,160)$ would be
reproduced under one of two readings; it was not, and the pair this
paper prints, in Note~\ref{note:notanerrorbar} and in the Summary, is
$(7.3,\,158)$, the measured one. C4 likewise changed a number this
paper prints: $1/(2\langle\log N\rangle)$ measures $0.036038$ against
the published $0.0358$, and it is the measured value that appears in
Measurement~\ref{meas:exponent} and in
Note~\ref{note:notanerrorbar}. L3 registered no shape at all but a
size --- that the floor would survive permutation within a tenth ---
and its refutation is what \S\ref{sec:placebo} was written around:
Note~\ref{note:floorsignal}, Lemma~\ref{lem:permfloor} and
Note~\ref{note:collapse} all follow from it.

\end{document}
